% RAFTcorr User Guide -- preamble (pdflatex, xelatex or lualatex).

\documentclass[11pt,letterpaper]{article}

% ---------------------------------------------------------------- fonts
\usepackage[T1]{fontenc}
\usepackage[utf8]{inputenc}
\usepackage{lmodern}
\usepackage{amssymb}
\usepackage{amsmath}
\usepackage{textcomp}
\usepackage{microtype}
\DisableLigatures[-]{encoding=T1,family=tt*}   % --roi stays two hyphens in typewriter text

% ---------------------------------------------------------------- layout
\usepackage[margin=1in,headheight=14pt]{geometry}
\usepackage{fancyhdr}
\pagestyle{fancy}
\fancyhf{}
\fancyhead[L]{\small RAFTcorr User Guide}
\fancyhead[R]{\small\thepage}
\renewcommand{\headrulewidth}{0.4pt}
\setlength{\parskip}{4pt}

% ---------------------------------------------------------------- graphics
\usepackage{graphicx}
\graphicspath{{figures/}}
\usepackage{float}
\usepackage[font=small,labelfont=bf,margin=0.5in]{caption}

% ---------------------------------------------------------------- color
\usepackage{xcolor}
\definecolor{raftaccent}{RGB}{79,70,229}     % the application's accent (light theme)
\definecolor{raftgreen}{RGB}{22,163,74}
\definecolor{raftred}{RGB}{220,38,38}
\definecolor{raftamber}{RGB}{180,83,9}
\definecolor{codebg}{RGB}{245,245,247}
\definecolor{keybg}{RGB}{243,244,246}

% ---------------------------------------------------------------- hyperlinks
\usepackage[colorlinks=true,linkcolor=raftaccent,urlcolor=raftaccent,
            citecolor=raftaccent]{hyperref}

% ---------------------------------------------------------------- callouts
\usepackage[most]{tcolorbox}
\newtcolorbox{tipbox}{
  colback=raftgreen!7!white, colframe=raftgreen, boxrule=0.5pt,
  arc=1pt, left=4pt, right=4pt, top=3pt, bottom=3pt,
  title={\textbf{Tip}}, fonttitle=\bfseries\small}
\newtcolorbox{warnbox}{
  colback=raftred!7!white, colframe=raftred, boxrule=0.5pt,
  arc=1pt, left=4pt, right=4pt, top=3pt, bottom=3pt,
  title={\textbf{Warning}}, fonttitle=\bfseries\small}
\newtcolorbox{notebox}{
  colback=raftaccent!7!white, colframe=raftaccent, boxrule=0.5pt,
  arc=1pt, left=4pt, right=4pt, top=3pt, bottom=3pt,
  title={\textbf{Note}}, fonttitle=\bfseries\small}

% ---------------------------------------------------------------- code
\usepackage{listings}
\lstset{
  basicstyle=\ttfamily\small,
  backgroundcolor=\color{codebg},
  frame=single, framerule=0pt,
  breaklines=true, breakindent=0pt,
  columns=fullflexible, keepspaces=true, showstringspaces=false,
  xleftmargin=4pt, xrightmargin=4pt,
  aboveskip=6pt, belowskip=6pt,
  literate={~}{{\textasciitilde}}1,
}

% ---------------------------------------------------------------- misc
\usepackage{multicol}
\usepackage{enumitem}
\setlist{topsep=3pt,itemsep=2pt,parsep=0pt}
\usepackage{booktabs}
\usepackage{tabularx}
\usepackage{array}

% ---------------------------------------------------------------- interface vocabulary
% \ui{Run DIC Analysis}     a label, button or field exactly as the application shows it
% \menu{File}{Save Project} a menu path (any number of levels: \menu{Help}{Models…} uses two)
% \kbd{Ctrl+S}              a key or a key combination
% \file{results.h5}         a file name, folder or command-line word
% \screenshot[width]{file}{caption}{label}   a figure from figures/
\newcommand{\ui}[1]{\textbf{#1}}
\newcommand{\menu}[2]{\textbf{#1}\,$\triangleright$\,\textbf{#2}}
\newcommand{\menuthree}[3]{\textbf{#1}\,$\triangleright$\,\textbf{#2}\,$\triangleright$\,\textbf{#3}}
\newcommand{\kbd}[1]{\tcbox[on line, boxsep=0pt, left=3pt, right=3pt, top=1.5pt, bottom=1.5pt,
  colback=keybg, colframe=black!30, boxrule=0.4pt, arc=2pt]{\footnotesize\ttfamily #1}}
\newcommand{\file}[1]{\texttt{#1}}
\newcommand{\screenshot}[4][\textwidth]{%
  \begin{figure}[htbp]
    \centering
    \includegraphics[width=#1]{#2}
    \caption{#3}
    \label{#4}
  \end{figure}}

% Unicode characters used in labels of the application
\DeclareUnicodeCharacter{2026}{\ldots}
\DeclareUnicodeCharacter{03B5}{\ensuremath{\varepsilon}}
\DeclareUnicodeCharacter{03C3}{\ensuremath{\sigma}}
\DeclareUnicodeCharacter{03C9}{\ensuremath{\omega}}
\DeclareUnicodeCharacter{0394}{\ensuremath{\Delta}}
\DeclareUnicodeCharacter{2081}{\ensuremath{_1}}
\DeclareUnicodeCharacter{2082}{\ensuremath{_2}}
\DeclareUnicodeCharacter{2080}{\ensuremath{_0}}
\DeclareUnicodeCharacter{2192}{\ensuremath{\rightarrow}}
\DeclareUnicodeCharacter{00D7}{\ensuremath{\times}}
\DeclareUnicodeCharacter{2264}{\ensuremath{\leq}}
\DeclareUnicodeCharacter{2265}{\ensuremath{\geq}}

\newcommand{\raftversion}{2.0.1}
   % \raftversion, written by build.py from src/raftcorr/_version.py

\hypersetup{pdftitle={RAFTcorr User Guide}, pdfauthor={Zixiang Tong}}

\begin{document}

% ---------------------------------------------------------------- title page
\begin{titlepage}
  \centering
  \vspace*{1.2in}
  \includegraphics[width=1.35in]{raftcorr-icon.png}\par
  \vspace{0.35in}
  {\fontsize{40}{46}\selectfont\bfseries RAFTcorr\par}
  \vspace{0.12in}
  {\Large\color{raftaccent} Learned Digital Image Correlation\par}
  \vspace{0.08in}
  {\large User Guide for the desktop application, the command line and the Python library\par}
  \vspace{0.5in}
  {\large Version \raftversion\par}
  \vfill
  {\large Zixiang Tong\par}
  {\href{mailto:zachtong@utexas.edu}{zachtong@utexas.edu}\par}
  \vspace{0.06in}
  {Dr.\ Jin Yang Group, The University of Texas at Austin\par}
  \vspace{0.12in}
  {\href{https://github.com/zachtong/RAFTcorr}{github.com/zachtong/RAFTcorr}\par}
  \vspace*{0.4in}
\end{titlepage}

\begin{abstract}
\noindent
This guide follows a complete RAFTcorr analysis, from installing the application to exporting
displacement and strain. Each section maps to a panel, a window or a menu of the application, so
it can be read with RAFTcorr open beside it. The two-page \emph{Quick Guide} holds the same steps
in condensed form; the last sections cover the command line, the Python library and
troubleshooting.
\end{abstract}

{\small
\begin{multicols}{2}
\tableofcontents
\end{multicols}}
\newpage

\section{Overview}
\label{sec:overview}

RAFTcorr measures full-field displacement and strain from a sequence of images of a deforming
specimen. The first image of the sequence is the reference image. Every other image is compared
with it, or with a reference frame earlier in the sequence, by a learned correlation model: a neural network trained
on speckle images whose displacement is known exactly. The model returns one displacement vector
for every pixel of the region of interest, so there is no subset size and no mesh to choose.
Strain follows from the displacement with a virtual strain gauge, in the strain window.

The same engine runs behind three interfaces: the desktop application, which most of this guide
follows; the \file{raftcorr} command line; and the Python library (Section~\ref{sec:cli}). A
project saved by one of them opens in the others.

\begin{notebox}
RAFTcorr is two-dimensional digital image correlation: one camera, a planar specimen, motion in
the image plane. It does not reconstruct out-of-plane motion.
\end{notebox}

\subsection{The four models}
\label{sec:overview-models}

RAFTcorr ships four models, \file{RAFTcorr-s1}, \file{-s2}, \file{-s4} and \file{-s8}. They share
one network architecture and differ in one number, the encoder downsampling factor $s$: the
network describes both images on a grid $s$ times coarser than the pixels, and the correlation
search at its coarsest level looks $32\,s$ px away. That distance is the model's nominal
\emph{reach}.

\begin{table}[htbp]
  \centering
  \caption{The four models. The \ui{Model} list shows each with its reach, for example
  \ui{RAFTcorr-s4 (128 px)}.}
  \label{tab:models}
  \begin{tabular}{lccl}
    \toprule
    Model & $s$ & Reach ($32\,s$) & Use it for \\
    \midrule
    RAFTcorr-s1 & 1 & 32 px  & small displacements, finest detail (strain concentrations) \\
    RAFTcorr-s2 & 2 & 64 px  & small to moderate displacements, fine detail, faster than s1 \\
    RAFTcorr-s4 & 4 & 128 px & a balance of detail and reach (the default) \\
    RAFTcorr-s8 & 8 & 256 px & large motion and rotation, or no prior knowledge \\
    \bottomrule
  \end{tabular}
\end{table}

A smaller $s$ resolves shorter wavelengths of the displacement; a larger $s$ reaches further. The
reach has to cover the largest displacement of any image pair, so the choice depends on the
tracking mode too (Section~\ref{sec:workflow-model}). When you do not know how far the specimen
moves, \ui{Suggest a Model} measures the motion of your sequence with s8 and picks the finest
model whose reach covers it.

Each model file is about 20\,MB. RAFTcorr downloads a model the first time an analysis uses it
and keeps it for later runs (Section~\ref{sec:installing-models}).

\subsection{An analysis in eight steps}
\label{sec:overview-steps}

The main window has three columns: settings on the left, the image in the middle, the run controls
and the display settings on the right. An analysis goes through it from left to right.

\begin{enumerate}
  \item \textbf{Load the images} in the \ui{IMAGES} section: a folder of images, sorted in natural
    order. The first image is the reference image (Section~\ref{sec:images}).
  \item \textbf{Choose the tracking mode and the model} in \ui{WORKFLOW TYPE}. Do this before the
    region of interest: the tracking mode decides which frames are reference frames
    (Section~\ref{sec:workflow}).
  \item \textbf{Define the region of interest} on frame~1 in \ui{REGION OF INTEREST}, and on
    later frames that must leave a region out, such as a crack or a hole that opens
    (Section~\ref{sec:roi}).
  \item \textbf{Check the parameters} in \ui{PARAMETERS} and \ui{ADVANCED}. The defaults suit most
    sequences (Section~\ref{sec:parameters}).
  \item \textbf{Run} with \ui{Run DIC Analysis}, and follow the run in \ui{PROGRESS} and \ui{LOG}
    (Section~\ref{sec:running}).
  \item \textbf{Review the displacement} on the canvas: each component, the magnitude, and the
    residual, which shows how well each frame was tracked (Sections~\ref{sec:canvas}
    and~\ref{sec:results}).
  \item \textbf{Compute the strain} in the strain window, which opens when the run finishes, and
    measure points, lines and areas in its \ui{Analysis} tab (Sections~\ref{sec:strain}
    and~\ref{sec:analysis}).
  \item \textbf{Export and save}: data files, figures, animations and a report with
    \ui{Export Results}, and the whole session with \menu{File}{Save Project}
    (Sections~\ref{sec:export} and~\ref{sec:projects}).
\end{enumerate}

\screenshot{fig-main-window.png}{The main window after a run, showing the vertical displacement on
the last frame. The numbers mark the five regions described in the text: (1) the \ui{IMAGES}
section; (2) the settings column; (3) the canvas with the frame navigator; (4) the run controls and
\ui{PROGRESS}; (5) \ui{FIELD}, \ui{VISUALIZATION}, \ui{PHYSICAL UNITS} and \ui{LOG}.}{fig:main-window}

The five regions of Figure~\ref{fig:main-window}:

\begin{description}[leftmargin=0.6cm,style=nextline]
  \item[1. IMAGES]
    The drop zone (\ui{Drop image folder or Browse}), the \ui{Natural Sort (1, 2, …, 10)}
    checkbox and the image list. The list numbers the frames from 1 and has a \ui{Region}
    column, whose button adds or edits the region of interest of each frame. See
    Section~\ref{sec:images}.
  \item[2. The settings column]
    Four sections that fold open and closed: \ui{WORKFLOW TYPE} (tracking mode, model,
    \ui{Suggest a Model}), \ui{REGION OF INTEREST} (drawing tools and mask files),
    \ui{PARAMETERS} (iterations, smoothing, device) and \ui{ADVANCED} (margin, GPU memory, tile
    blending, correlation; folded by default). See Sections~\ref{sec:workflow},
    \ref{sec:roi} and~\ref{sec:parameters}.
  \item[3. The canvas]
    The current frame with the region of interest or the displacement field drawn over it, the
    zoom buttons and a read-out of the position and value under the cursor at the top, and the
    frame navigator at the bottom. See Section~\ref{sec:canvas}.
  \item[4. Run controls and PROGRESS]
    \ui{Run DIC Analysis}, \ui{Cancel}, \ui{Export Results} and \ui{Open Strain Window}, and below
    them the progress of a run. See Sections~\ref{sec:running}, \ref{sec:strain}
    and~\ref{sec:export}.
  \item[5. Display and log]
    \ui{FIELD} chooses what is drawn (\ui{Disp U}, \ui{Disp V}, \ui{Magnitude}, \ui{Residual}) and
    where (\ui{Show on}); \ui{VISUALIZATION} holds the colormap, the colour range and the opacity;
    \ui{PHYSICAL UNITS} converts pixels and frames to lengths and seconds; \ui{LOG} lists every
    message of the session with its time, warnings in yellow and errors in red. See Section~\ref{sec:results}.
\end{description}

The menu bar has three menus: \ui{File} (images, projects, export), \ui{Settings} (theme and
language) and \ui{Help} (models, self-test, version). Section~\ref{sec:settings} lists them with
their keyboard shortcuts.

\subsection{Conventions}
\label{sec:conventions}

These rules hold everywhere: on the screen, in exported files, on the command line and in
Python.

\begin{itemize}
  \item \textbf{Displacement} $(u, v)$ goes from the reference image to the deformed image: a
    material point at pixel $X$ in the reference image appears at $X + (u, v)$ in the deformed
    image. $+u$ points right and $+v$ points down, as the image rows and columns run. The unit is
    the pixel, unless you set physical units (Section~\ref{sec:physical-units}).
  \item \textbf{Lagrangian storage.} Every field is stored in the coordinates of the reference
    image: the value at a pixel belongs to the material point that sits at that pixel in frame~1.
    Incremental tracking composes its steps into these coordinates as well. Choosing
    \ui{Deformed frame} in \ui{Show on} moves the field to the current positions of the points for
    display only; the stored data do not change.
  \item \textbf{Frame numbers.} The application counts frames from 1: the image list, the frame
    navigator and the log call the reference image frame~1. File names written by
    the export dialog count the same way (\file{frame\_0001} is the reference image). Inside the
    arrays of the exported data files, and in Python, index $k$ is image $k$ counted from 0, so
    frame~$n$ on the screen is index $n-1$. The command line uses both: its progress lines, its
    messages and its \file{suggest} report count from 1, its frame option \file{-{}-keys} from 0
    (Section~\ref{sec:cli}).
  \item \textbf{Shear strain} $\varepsilon_{xy}$ is the tensor component, half the engineering
    shear strain $\gamma_{xy}$ (Section~\ref{sec:strain}).
  \item \textbf{Grey levels.} RAFTcorr compares grey levels from 0 to 255. Colour images are
    converted to luminance, and images with more than 8 bits are scaled with one factor for the
    whole sequence, taken from the reference image (Section~\ref{sec:images}).
\end{itemize}

\subsection{Where RAFTcorr keeps its files}
\label{sec:overview-files}

RAFTcorr writes into the image folder only what you save there yourself. Its own files are in
your user account, so no administrator rights are needed.

\begin{description}[leftmargin=0.6cm,style=nextline]
  \item[Model files]
    \file{\%LOCALAPPDATA\%\textbackslash RAFTcorr\textbackslash models} on Windows,
    \file{\textasciitilde/.raftcorr/models} on macOS and Linux. \menu{Help}{Models…} shows the
    folder and opens it (Section~\ref{sec:installing-models}).
  \item[Settings]
    The theme, the language and the last export folder, in the settings store of your user
    account: the registry key
    \file{HKEY\_CURRENT\_USER\textbackslash Software\textbackslash RAFTcorr} on Windows, a property
    list in \file{\textasciitilde/Library/Preferences} on macOS.
  \item[Log file]
    \file{RAFTcorr.log}, however RAFTcorr was installed, in this folder:\newline
    Windows: \file{\%LOCALAPPDATA\%\textbackslash RAFTcorr\textbackslash logs}\newline
    macOS: \file{\textasciitilde/Library/Application Support/RAFTcorr/logs}\newline
    Linux: \file{\$XDG\_DATA\_HOME/RAFTcorr/logs}, by default
    \file{\textasciitilde/.local/share/RAFTcorr/logs}\newline
    At 2\,MB the
    file is renamed and a new one started; the three most recent older files are kept
    (Section~\ref{sec:troubleshooting}).
  \item[Windows installation]
    \file{\%LOCALAPPDATA\%\textbackslash Programs\textbackslash RAFTcorr} by default, with the
    private Python and \file{install.log}, the log of the setup (Section~\ref{sec:installing}).
  \item[Large results]
    When the results of a run (displacement and residual) need more than 2\,GB of memory,
    RAFTcorr keeps them in memory-mapped files in a temporary folder named
    \file{raftcorr\_results\_\ldots}, which the \ui{LOG} names, and deletes the folder when you
    quit. Save a project or export the results to keep them (Sections~\ref{sec:projects}
    and~\ref{sec:export}).
\end{description}

\section{Installing RAFTcorr}
\label{sec:installing}

Pick the route that fits your computer. Every route gives you the same application, the same
\file{raftcorr} command line and the same model files.

\begin{center}
\begin{tabularx}{\linewidth}{@{}p{0.4\linewidth}X@{}}
  \toprule
  You have & Use \\
  \midrule
  Windows 10 (version 1809 or later) or 11, 64-bit & the Setup program, \file{RAFTcorr-Windows-Setup.exe}
    (Section~\ref{sec:installing-windows}) \\
  A Mac with Apple silicon, macOS 14 or later & the disk image, \file{RAFTcorr-macOS.dmg}
    (Section~\ref{sec:installing-macos}) \\
  Python 3.10 or later, on any system & \file{pip install raftcorr}
    (Section~\ref{sec:installing-pip}) \\
  A copy of the source code & \file{install.bat} on Windows, or a conda environment
    (Section~\ref{sec:installing-source}) \\
  \bottomrule
\end{tabularx}
\end{center}

Both downloads are on the release page, \url{https://github.com/zachtong/RAFTcorr/releases/latest}.

\subsection{Windows: the Setup program}
\label{sec:installing-windows}

The Setup program is small: it holds a private Python and RAFTcorr itself, but not PyTorch, the
library that runs the models. Which PyTorch build fits depends on your graphics card, so Setup
chooses and downloads it while it installs. It needs an internet connection once, for PyTorch,
the other Python packages and the model files.

\begin{enumerate}
  \item Download \file{RAFTcorr-Windows-Setup.exe} from the release page and run it. It installs
    for your own user account and asks for no administrator password, so it works on a lab
    computer too. Windows may warn you first; see the note below.
  \item Accept the licence and keep the proposed folder,
    \file{\%LOCALAPPDATA\%\textbackslash Programs\textbackslash RAFTcorr}, or choose another one.
  \item Choose the additional tasks. \ui{Add the "raftcorr" command to PATH (new terminals)} is
    ticked by default; \ui{Create a desktop shortcut} is not. Click \ui{Install}.
  \item Setup copies its files, then opens a console window and installs the Python components in
    three steps: \file{[1/3] PyTorch}, \file{[2/3] RAFTcorr} and its dependencies,
    \file{[3/3] model files}. The status line reads \emph{``Installing PyTorch, RAFTcorr and the
    model files (up to 3 GB with a GPU; this takes a while)...''}. Wait for the console window to
    close by itself.
  \item On the last page, keep \ui{Launch RAFTcorr} ticked to start the application. The box is
    offered only when the Python components were installed.
\end{enumerate}

\begin{warnbox}
Do not close the console window while it works: that stops the installation of the Python
components. If it happened, run \ui{Repair RAFTcorr} from the Start menu.
\end{warnbox}

\subsubsection*{Which PyTorch}
Setup asks the NVIDIA driver (the program \file{nvidia-smi}) which CUDA version it supports and
installs the matching PyTorch build:

\begin{center}
\begin{tabular}{ll}
  \toprule
  The driver reports & PyTorch build \\
  \midrule
  CUDA 12.8 or later, or a GeForce RTX 50 series card & CUDA 12.8 \\
  CUDA 12.0 to 12.7 & CUDA 12.6 \\
  CUDA 11.x & CUDA 11.8 \\
  CUDA 10 or older, no CUDA version, or no NVIDIA driver & CPU build \\
  \bottomrule
\end{tabular}
\end{center}

The console line of step \file{[1/3]} says which case applied, for example
\emph{``[1/3] PyTorch: driver supports CUDA 12.4: PyTorch with CUDA 12.6''}. Without an NVIDIA GPU, RAFTcorr runs
on the CPU, more slowly. If the driver is too old, the line ends with \emph{``(update the NVIDIA
driver for the GPU)''}: update the driver, then run \ui{Repair RAFTcorr}.

\subsubsection*{What you get}
\begin{itemize}
  \item \ui{RAFTcorr} in the Start menu, and on the desktop if you asked for it.
  \item \ui{Repair RAFTcorr} in the Start menu. It runs the second half of the setup again:
    detects the driver, installs PyTorch, RAFTcorr and the model files. A PyTorch of the build
    that fits the driver is kept, so nothing is downloaded for it again; a PyTorch of another
    build is uninstalled first and the right one installed, for example the CUDA build in place
    of the CPU build after a driver update. Use it when the setup was interrupted (no connection,
    console window closed) or after you installed or updated an NVIDIA driver.
  \item The \file{raftcorr} command in every new terminal, when the PATH task was ticked. Only the
    folder \file{bin} of the installation goes on the PATH; the private Python and its
    \file{pip} stay off it.
  \item \file{.raftcorr} project files open in RAFTcorr when you double-click them.
  \item \file{install.log} in the installation folder: the full output of every setup and
    repair, the PyTorch build chosen and, at the end, the versions found
    (\file{raftcorr info}, Section~\ref{sec:installing-check}).
\end{itemize}

If the Python components could not be installed, Setup says so: \emph{``RAFTcorr could not install
its Python components (code~N). Check the internet connection, then run `Repair RAFTcorr' from the
Start menu.''}, followed by the path of \file{install.log}. A model file that could not be
downloaded does not stop the setup; RAFTcorr fetches the model when it is first used.

\begin{notebox}
RAFTcorr is not code-signed, so Windows SmartScreen may stop the Setup program the first time
with \emph{``Windows protected your PC''}. Click \ui{More info}, then \ui{Run anyway}. SmartScreen
reports that the file is not yet widely downloaded; it is not a virus detection.
\end{notebox}

\subsubsection*{Upgrading}
Run the Setup program of the new version; there is no need to uninstall first. It installs into
the same folder, replaces the private Python with a clean one and installs PyTorch again, so it
needs the internet connection again. Your model files, settings and projects stay. Close RAFTcorr
before you start.

\subsubsection*{Uninstalling}
Open the Windows settings, \ui{Apps}, then \ui{Installed apps} (Windows~11) or
\ui{Apps \& features} (Windows~10), and uninstall \ui{RAFTcorr}. This removes the installation
folder with its private Python and \file{install.log}, the PATH entry and the \file{.raftcorr}
file association. The model files in
\file{\%LOCALAPPDATA\%\textbackslash RAFTcorr\textbackslash models} and your settings stay;
delete the folder by hand if you no longer need the models.

\subsection{macOS: the disk image}
\label{sec:installing-macos}

\begin{enumerate}
  \item Download \file{RAFTcorr-macOS.dmg} from the release page.
  \item Open it and drag \ui{RAFTcorr} onto \ui{Applications}.
  \item Start RAFTcorr from Applications, Launchpad or Spotlight.
\end{enumerate}

It needs a Mac with Apple silicon and macOS 14 (Sonoma) or later; Intel Macs are not supported.
The application is self-contained: Python, PyTorch and Qt are inside, and nothing else is
installed. PyTorch runs on the CPU, since a Mac has no CUDA. The model files are not inside; they
are downloaded the first time each is used (Section~\ref{sec:installing-models}). To upgrade,
drag the new version onto \ui{Applications} and replace the old one.

\begin{notebox}
The application is not notarised by Apple, so the first time macOS refuses to open it. Open
\ui{System Settings}, then \ui{Privacy \& Security}, scroll to the line about RAFTcorr, click
\ui{Open Anyway} and confirm. From then on it opens like any other application. A new version may
ask again.
\end{notebox}

The command line is inside the application. To use it as \file{raftcorr}, add an alias to your
shell, for example in \file{\textasciitilde/.zshrc}:

\begin{lstlisting}
alias raftcorr=/Applications/RAFTcorr.app/Contents/MacOS/RAFTcorr
\end{lstlisting}

Without arguments this command opens the application; with a command such as \file{run} or
\file{models} it is the command line of Section~\ref{sec:cli}. RAFTcorr does not register
\file{.raftcorr} project files on macOS: open a project with \menu{File}{Open Project…}.

\subsection{With pip}
\label{sec:installing-pip}

Use this route on Linux, or when you work in Python already. It needs Python 3.10 or later.
Install RAFTcorr into a virtual environment of its own, so its dependencies cannot collide with
other packages.

On a computer with an NVIDIA GPU, install PyTorch from the PyTorch index first, choosing the CUDA
build that your driver supports (the table of Section~\ref{sec:installing-windows};
\url{https://pytorch.org/get-started/locally/} lists the builds). The PyTorch on PyPI is the CPU
build on Windows, and \file{pip install raftcorr} alone would install it. On a Mac, skip this
step.

\begin{lstlisting}
pip install torch --index-url https://download.pytorch.org/whl/cu126
pip install raftcorr
raftcorr                     # start the application
raftcorr models --download   # fetch the model files now (otherwise on first use)
raftcorr selftest            # run every installed model on a synthetic pair
\end{lstlisting}

\file{python -m raftcorr} works like \file{raftcorr} when the script folder of pip is not on your
PATH. pip also installs \file{raftcorr-gui}, which starts the application without a console
window on Windows.

\begin{notebox}
On Windows, PySide6 and PyTorch install files up to about 150 characters deep below
\file{site-packages}. If pip stops with \emph{``No such file or directory''} and a hint about long
paths, create the environment in a shorter folder, for example
\file{C:\textbackslash venvs\textbackslash raftcorr}, or enable long paths in Windows
(\url{https://pip.pypa.io/warnings/enable-long-paths}).
\end{notebox}

If PyTorch is already installed as the CPU build, \file{pip install torch -{}-index-url ...} keeps
it. Remove it first with \file{pip uninstall torch}, then install the CUDA build.

\subsection{From the source code}
\label{sec:installing-source}

Clone the repository, \url{https://github.com/zachtong/RAFTcorr}. On Windows with Anaconda or
Miniconda, run \file{install.bat} in its folder. It creates the conda environment
\file{raftcorr-app} (Python 3.11), detects the GPU as the Setup program does, installs the
matching PyTorch build (a download of about 2.5\,GB), RAFTcorr in editable mode and the model
files, and writes \file{install.log} in the same folder. Like Repair, it keeps a PyTorch of the
matching build, so running it again downloads no PyTorch, and uninstalls a PyTorch of another
build before it installs the matching one. Start the application with
\file{run.bat}.

On any platform:

\begin{lstlisting}
conda env create -f environment.yml
conda activate raftcorr-app
pip install torch --index-url https://download.pytorch.org/whl/cu126  # or /cpu
pip install -e ".[dev]"
\end{lstlisting}

In a source checkout, RAFTcorr also looks for model files in its folder \file{models}.

\subsection{First launch}
\label{sec:installing-first}

Start RAFTcorr from the Start menu, Applications, \file{run.bat} or the \file{raftcorr} command.
The main window opens empty: the image list shows the drop zone \ui{Drop image folder or Browse},
the canvas is blank and the progress area reads \ui{Ready}. The window does not shrink below
1120 $\times$ 760 pixels. The first time, the interface follows the language of the operating
system when RAFTcorr has a translation for it; change it in \menu{Settings}{Language}
(Section~\ref{sec:settings}).

\subsection{The models}
\label{sec:installing-models}

You need not install the models by hand. The first time an analysis uses a model that is not on
the computer, RAFTcorr downloads it from the release \file{weights-v1} of the RAFTcorr repository.
The \ui{PROGRESS} area says so, for example \emph{``Downloading model RAFTcorr-s4: 5.2 of 20.7
MB''}, and the bar follows the download; then the bar starts again from 0 for the tracking.
\ui{Cancel} stops the download and leaves no partial file. Every file is checked against its
SHA-256 before it is used, and kept for later runs; the \ui{LOG} names the address and the file
it is saved to. \ui{Suggest a Model} downloads s8 the same way when it is missing, and shows the download in
its progress dialog.

\screenshot[0.75\textwidth]{fig-models-window.png}{The Models window
(\menu{Help}{Models…}). One row per model with its downsampling factor, its reach, its state and
a \ui{Download} button while it is missing; below, the model folder, the link to the release page
and the buttons \ui{Download All}, \ui{Import Model Files…}, \ui{Open Model Folder} and
\ui{Close}.}{fig:models-window}

\menu{Help}{Models…} opens the Models window (Figure~\ref{fig:models-window}):

\begin{description}[leftmargin=0.6cm,style=nextline]
  \item[The model rows]
    Each row shows the model, \emph{``s = 4, reach 128 px''}, and its state: \ui{Installed} (the
    tooltip gives the path) or \ui{Not installed}, with the size of the download. A missing model
    has a \ui{Download} button.
  \item[Download All]
    Downloads every missing model, one after the other, and stops at the first failure.
    \ui{Cancel Download} appears while a download runs. Closing the window also cancels it.
  \item[Import Model Files…]
    Installs \file{.pth} files that you downloaded by hand. The SHA-256 of a file says which model
    it is, so a renamed download such as \file{RAFTcorr-s4 (1).pth} is recognised and stored under
    its release name; any other file is refused with \emph{``Not imported: \ldots\ Only the model
    files of the release can be imported.''}
  \item[Open Model Folder]
    Opens the per-user model folder, shown above the buttons as \ui{Model folder}.
\end{description}

When a run, or \ui{Suggest a Model}, needs a model that is neither on the computer nor
downloadable, the Models window opens instead of an error box. A card at the top names the model
and the reason, for example \emph{``RAFTcorr-s4 is needed, but it is not installed and could not
be downloaded.''} followed by \emph{``no connection to the download server (\ldots); check the
internet connection''}. Once the model is there, the card reads \emph{``RAFTcorr-s4 is installed
now. Close this window and start again.''} Then click \ui{Run DIC Analysis} again.

\subsubsection*{A computer without internet}

\begin{enumerate}
  \item On a computer with internet, open the release page of the models,
    \url{https://github.com/zachtong/RAFTcorr/releases/tag/weights-v1}; the Models window links
    to it. Download the files you need, from \file{RAFTcorr-s1.pth} to \file{RAFTcorr-s8.pth},
    about 20\,MB each.
  \item Copy them to the offline computer.
  \item Click \ui{Import Model Files…} in the Models window, or run
    \file{raftcorr models -{}-import FILE ...} (Section~\ref{sec:cli}).
\end{enumerate}

The Windows Setup program itself downloads PyTorch, so it cannot complete without internet. The
macOS disk image needs the connection only for the model files.

\subsubsection*{Where RAFTcorr looks for the model files}

For each model, RAFTcorr takes the first verified copy it finds, in this order:

\begin{enumerate}
  \item a file path that you give (\file{-{}-model} on the command line, or the Python library);
  \item the folder named by the environment variable \file{RAFTCORR\_WEIGHTS};
  \item the folder \file{models} of a source checkout;
  \item the per-user folder, \file{\%LOCALAPPDATA\%\textbackslash RAFTcorr\textbackslash models}
    on Windows, \file{\textasciitilde/.raftcorr/models} on macOS and Linux;
  \item a download into the per-user folder.
\end{enumerate}

A file whose SHA-256 does not match the release is skipped, never used, and the log says so:
\emph{``\ldots\ exists but its SHA-256 does not match the RAFTcorr-s4 release; skipping it''}.
\file{raftcorr models} lists the four models, where each one is installed and the folders
searched.

\subsection{Checking an installation}
\label{sec:installing-check}

\menu{Help}{Self-test} runs every installed model on a synthetic pair of $256 \times 256$ px
speckle images with a known translation, $u = +2.30$\,px and $v = -1.20$\,px, on the device set
in \ui{PARAMETERS}. The result opens in a dialog and goes to the log: one line per model with the
measured $u$ and $v$ and \file{ok} or \file{WRONG}, then \emph{``All installed models pass.''} or
\emph{``Some models FAILED.''} The self-test downloads nothing: a model that is not installed is
listed as \file{not installed}, and with no model installed at all the test fails.

From a terminal:

\begin{lstlisting}
raftcorr selftest              # the same test, on the device "auto"
raftcorr selftest --gui        # also open and close the main window offscreen
raftcorr info                  # versions, CUDA and the GPU
\end{lstlisting}

\file{raftcorr selftest} ends with exit code 0 when every check passes. \file{raftcorr info}
prints \file{CUDA available: True} or \file{False}, and with a GPU its name and free memory.
Section~\ref{sec:cli-other} describes both commands.

\subsection{GPU or CPU}
\label{sec:installing-gpu}

RAFTcorr runs on an NVIDIA GPU through CUDA, or on the CPU. For the GPU you need an NVIDIA card,
a current NVIDIA driver and a CUDA build of PyTorch. The Windows Setup program and
\file{install.bat} choose that build for you; with pip you install it yourself
(Section~\ref{sec:installing-pip}). On a Mac, RAFTcorr runs on the CPU. The \ui{Device} setting,
\ui{Automatic} by default, picks the GPU when PyTorch sees one (Section~\ref{sec:parameters});
\file{raftcorr info} tells you whether it does.

\section{Loading images}
\label{sec:images}

An analysis starts from a folder that holds the image sequence. The \ui{IMAGES} section at the
top of the left sidebar loads the folder, sets the order of the images and lists them. The first
image of the list is the \emph{reference image}: every displacement is measured from it and
reported in its coordinates. On screen it is frame~1; in exported data files and in Python it is
index~0.

\subsection{Opening a folder}

There are three ways to open a folder:
\begin{itemize}
  \item drag the folder onto the drop zone (\ui{Drop image folder or Browse}); if you drop a file
        instead, RAFTcorr opens the folder that contains it;
  \item click the drop zone and choose the folder in the \ui{Select Image Folder} dialog;
  \item choose \menu{File}{Open Image Folder…} (\kbd{Ctrl+I}; \kbd{Cmd+I} on macOS).
\end{itemize}
From a terminal, \file{raftcorr gui FOLDER} starts the application with that folder open
(Section~\ref{sec:cli}).

RAFTcorr reads every image file directly inside the folder. It does not look into subfolders,
and it skips hidden files (names that start with a dot). After loading, the badge next to the
\ui{IMAGES} title shows the number of images, the canvas shows frame~1 fitted to the view, and
the \ui{LOG} in the right sidebar reports the number of images, their size in pixels and the
grey-level scale (next subsection).

\begin{warnbox}
Opening a folder starts a new session. The regions of interest, the results, the probes, the
load data and the model suggestion of the previous images are discarded. Save a project first
if you want to keep them (Section~\ref{sec:projects}).
\end{warnbox}

\subsection{Formats, bit depths and colour}

RAFTcorr reads PNG (\file{.png}), TIFF (\file{.tif}, \file{.tiff}), BMP (\file{.bmp}) and JPEG
(\file{.jpg}, \file{.jpeg}) files; the extension is not case sensitive. Of a multi-page TIFF file
only the first page is read.

The network works on grey levels from 0 to 255. RAFTcorr converts every image with the same
rules, so that the brightness of a speckle does not change from one image to the next because
of the conversion:
\begin{itemize}
  \item \textbf{Colour images} become their luminance ($0.299\,R + 0.587\,G + 0.114\,B$, the
        ITU-R BT.601 weights). An alpha channel is dropped.
  \item \textbf{8-bit images} are used as they are.
  \item \textbf{Images stored with more bits} (16-bit PNG or TIFF files) are scaled by their
        \emph{effective} bit depth: the smallest of 8, 10, 12, 14 and 16 bits that holds the
        brightest pixel of the reference image. A 12-bit camera that writes 16-bit files gives
        values up to 4095, so RAFTcorr maps 4095 (not 65535) to 255 and the speckle keeps its
        contrast.
  \item \textbf{Floating-point images} with values from 0 to 1 are multiplied by 255; any other
        range is divided by the largest value of the reference image and multiplied by 255.
\end{itemize}
The scaled values stay floating-point numbers (they are not rounded back to 8 bits) and are
clipped to the range 0 to 255. The canvas shows the images after this conversion, in grey. The
log line after loading names the scale in use, for example \file{8-bit}, or
\file{12-bit (stored as uint16)} for the camera above.

\begin{notebox}
The scale is decided once, from the reference image, and applied to the whole sequence. If a
later image is brighter than the bit depth of the reference image allows (a 12-bit reference
followed by images with values above 4095), its brightest pixels are clipped to 255. Keep the
exposure constant over the sequence.
\end{notebox}

\subsection{Image size}

All images of a sequence must have the size of the reference image. If one does not, RAFTcorr
shows the \ui{Cannot open the images} dialog and loads nothing. The message names the reference
image with its size and up to five of the images that differ, each with its frame number and its
size, for example \emph{``every frame must have the size of frame 1, a1.png (80 x 64 px);
different: frame 2, a2.png (82 x 64 px)''}.

The models also need a minimum image size: each side must be at least $16\,s$~px, which is
16, 32, 64 and 128~px for RAFTcorr-s1, -s2, -s4 and -s8. A smaller image stops the run with an
error.

\subsection{Sort order}

The order of the list is the order of the sequence, and the first image becomes the reference
image. The \ui{Natural Sort (1, 2, …, 10)} checkbox, on by default, sorts by the numbers inside
the file names: \file{img2.png} comes before \file{img10.png}. Unticked, the list is in plain
alphabetical order, which puts \file{img10.png} before \file{img2.png}. For names with
zero-padded numbers (\file{img\_0002.png}, \file{img\_0010.png}) both orders are the same.

\begin{warnbox}
When changing \ui{Natural Sort} puts the images in a different order, RAFTcorr reloads them and,
like opening a folder, starts a new session: the regions of interest, the results and the probes
are discarded. If there is anything to lose (regions of interest, results, probes, load data or
a running analysis), it asks first with \ui{Change the image order?}; \ui{No} keeps everything
and sets the checkbox back. When the order stays the same, as with zero-padded names, the
checkbox changes nothing else. Set the order before you draw a region of interest.
\end{warnbox}

\subsection{The image list}

\screenshot[0.45\textwidth]{fig-images.png}{The \ui{IMAGES} section after opening a folder of 44
images: the count badge, the drop zone, the \ui{Natural Sort} checkbox and the image list with its
\ui{\#}, \ui{Filename} and \ui{Region} columns.}{fig:images}

The list below the drop zone (Figure~\ref{fig:images}) has three columns:
\begin{description}[leftmargin=1.2cm,style=nextline]
  \item[\#] The frame number, counted from 1 and shown with two digits (\file{01}, \file{02},
    \ldots), the same number as on the frame slider below the canvas.
  \item[Filename] The file name. The rows of the reference frames (the frames that later frames
    are compared with, Section~\ref{sec:workflow}) are tinted with the accent colour and their
    names are bold. Which frames these are depends on the tracking mode, and the highlight
    follows it when you change the mode.
  \item[Region] A button that shows whether the frame has a mask of its own:
    \begin{itemize}
      \item \textcolor{raftred}{\ui{Need}} (red): frame~1 has no region of interest yet. A run
            cannot start without one.
      \item \textcolor{raftgreen}{\ui{Edit}} (green): the frame has a mask of its own.
      \item \ui{Add} (grey): the frame has no mask of its own. This is the normal state of every
            frame after the first; a mask there is optional.
    \end{itemize}
    Clicking the button shows that frame on the canvas and opens its mask for editing
    (Section~\ref{sec:roi-editing}).
\end{description}

Clicking the number or the name of a row shows that frame on the canvas. When the list has the
keyboard focus, \kbd{Up}, \kbd{Down}, \kbd{Page Up}, \kbd{Page Down}, \kbd{Home} and \kbd{End}
move through the frames. \kbd{Ctrl}-click (\kbd{Cmd}-click on macOS) and \kbd{Shift}-click
select several rows.

Right-clicking the selected rows opens a menu:
\begin{description}[leftmargin=1.2cm,style=nextline]
  \item[Import Region of Interest for $n$ frame(s)] Opens a file dialog. Choose exactly as many
    mask files as rows are selected; with any other number nothing is imported and the log says
    that the counts must match. The selected frames, from the first to the last, receive the
    files in the natural order of their names (\file{mask2.png} before \file{mask10.png}),
    whatever order you clicked them in. Check the \ui{Region} buttons afterwards, or use \ui{Batch Import}, which shows every
    pairing before it applies them (Section~\ref{sec:roi}).
  \item[Clear Region of Interest] Removes the masks of the selected frames. The entry tells how
    many of them have one.
  \item[Delete $n$ image(s)] Removes the selected images from the sequence. The files stay on
    disk. The remaining frames are numbered again, their masks follow their images and the
    probes are kept; the results are discarded, because their frame numbers no longer match.
    Deleting frame~1 makes the next image the reference image; the region of interest of the
    deleted frame goes with it, so the new frame~1 keeps only a mask it had of its own.
\end{description}

\section{Workflow type: tracking mode and model}
\label{sec:workflow}

The \ui{WORKFLOW TYPE} section, below the image list (Figure~\ref{fig:workflow}), holds the two
choices that shape a run:
which image each frame is compared with (the \ui{Tracking Mode}) and which network compares them
(the \ui{Model}). It sits above \ui{REGION OF INTEREST} because the tracking mode decides which
frames are reference frames, and the hint of the region of interest follows it
(Section~\ref{sec:roi}). The card in the top-left corner of the canvas repeats both choices
(Section~\ref{sec:canvas}).

\screenshot[0.45\textwidth]{fig-workflow.png}{The \ui{WORKFLOW TYPE} section: the \ui{Tracking
Mode} and \ui{Model} lists and the \ui{Suggest a Model} button.}{fig:workflow}

\subsection{Tracking Mode}

A sequence of $n$ images is processed as $n-1$ image pairs, one for every frame after the first.
Each pair holds a \emph{reference frame} $r$ and a later frame $k$; the network measures the
displacement from $r$ to $k$. The tracking mode decides which reference frame each frame gets.
Every mode computes one pair per frame, and every mode reports its results in the coordinates of
frame~1. Table~\ref{tab:tracking-modes} summarises the choices.

\begin{table}[htbp]
  \centering
  \caption{What each frame $k$ is compared with. Frame numbers as in the image list.}
  \label{tab:tracking-modes}
  \begin{tabularx}{\textwidth}{>{\raggedright\arraybackslash}p{3.6cm}
                               >{\raggedright\arraybackslash}X
                               >{\raggedright\arraybackslash}p{3.4cm}}
    \toprule
    Tracking Mode & Frame $k$ is compared with & Reference frames \\
    \midrule
    \ui{Accumulative} & frame 1 & 1 \\
    \ui{Incremental}, \ui{Every Frame} & frame $k-1$ & every frame but the last \\
    \ui{Incremental}, \ui{Every N Frames} & frame $1 + N\lfloor (k-2)/N \rfloor$:
      frames 2 to $N+1$ with frame 1, frames $N+2$ to $2N+1$ with frame $N+1$, \ldots
      & $1, N+1, 2N+1, \ldots$ \\
    \ui{Incremental}, \ui{Custom Frames} & the last reference frame before $k$
      & 1 and the frames you list \\
    \bottomrule
  \end{tabularx}
\end{table}

\subsubsection*{Accumulative (default)}

Every frame is compared directly with frame~1. Nothing is composed, so an error in one frame
does not pass on to the next. This is the mode to use whenever the displacement from frame~1
stays within the reach of the model (Section~\ref{sec:workflow-model}) in every frame of the
sequence.

\subsubsection*{Incremental}

Each frame is compared with a reference frame nearer to it than frame~1, and the steps are
composed back to frame~1. Choose it when the motion from frame~1 grows beyond the reach of the model, as in a
large stretch or a large rotation. The price is that every step adds its own error to all the
frames measured after it.

Choosing \ui{Incremental} shows the \ui{Reference Update} list, which sets how often the
reference frame changes:
\begin{description}[leftmargin=1.2cm,style=nextline]
  \item[Every Frame (default)]
    Each frame is compared with the one before it. Each step is the motion between two
    consecutive images, the shortest step the sequence offers.
  \item[Every N Frames]
    Shows the \ui{Interval} box (2 to 100 frames, default 5). The reference frame changes every
    $N$ frames: with $N = 5$, frames 2 to 6 are compared with frame~1, frames 7 to 11 with
    frame~6, and so on. Fewer steps are composed than with \ui{Every Frame}, but each step is
    longer. Choose $N$ so that the motion over $N$ frames stays within the reach of the model.
  \item[Custom Frames]
    Shows the \ui{Reference Frames} field. Type the frame numbers, as in the image list,
    separated by commas (for example \file{5, 10, 20}); the setting applies when you press
    \kbd{Enter} or leave the field. Frame~1 is always the first reference frame, and each frame
    is compared with the last reference frame before it. Use it when the motion is uneven: put
    the reference frames closer together where the specimen moves fast. The last frame cannot be
    a reference frame, because no frame follows it, so with $n$ images the numbers must lie
    between 1 and $n-1$. A list that breaks this rule, or that is not made of whole numbers, is
    refused: the field gets a red border, its tooltip and a warning in the log say what is wrong,
    for example \emph{``Invalid reference frames 44: use frames 1 to 43, separated by commas (the
    last frame cannot be a reference).''}. While the field holds refused text, \ui{Run DIC
    Analysis} does not start: a \ui{Cannot start analysis} dialog repeats the reason. Reference
    frames
    that the loaded images cannot use, because a project or an earlier folder had more images,
    are dropped, and the log names them.
\end{description}

\subsubsection*{How the steps are composed}

For a pair $(r, k)$ the network returns the displacement $\mathbf{d}_{rk}$ from frame $r$ to
frame $k$, in the coordinates of frame $r$. A point $\mathbf{X}$ of the region of interest of
frame~1 sits at $\mathbf{X} + \mathbf{U}_r(\mathbf{X})$ in frame $r$, so its displacement to
frame $k$ is
\begin{equation*}
  \mathbf{U}_k(\mathbf{X}) = \mathbf{U}_r(\mathbf{X})
      + \mathbf{d}_{rk}\bigl(\mathbf{X} + \mathbf{U}_r(\mathbf{X})\bigr),
\end{equation*}
with $\mathbf{d}_{rk}$ interpolated bilinearly between pixels. In every mode the result is the
displacement of the points of frame~1 (Lagrangian coordinates): $+u$ to the right, $+v$
downward, in pixels.

Composition has three consequences worth knowing:
\begin{itemize}
  \item The region of interest of frame~1 is carried to each reference frame with the
        displacement measured so far, and the next pairs are solved over that carried region.
        A mask you give a reference frame replaces the carried region there
        (Section~\ref{sec:roi-later-frames}). At the start of a run, the \ui{LOG} lists the
        reference frames that carry the region of frame~1: \emph{Reference frames without a
        Region of Interest of their own carry the one of frame 1 forward}, followed by their
        numbers.
  \item A point that is lost in a reference frame (it left the image, or a mask excludes it)
        stays empty in every frame measured from that reference frame.
  \item If a reference frame has no result, the frames compared with it are skipped, and the
        log says so.
\end{itemize}

\begin{notebox}
After a run, RAFTcorr watches the residual (the grey-level mismatch left after the
displacement maps one image onto the other) and warns when it jumps from one frame on. In
accumulative tracking the warning suggests incremental tracking or a model with a larger $s$;
in incremental tracking it suggests more reference frames or a model with a larger $s$
(Section~\ref{sec:tracking-quality}).
\end{notebox}

\subsection{Model}
\label{sec:workflow-model}

The \ui{Model} list holds the four models of Table~\ref{tab:models}, each with its reach in
parentheses, as \ui{RAFTcorr-s4 (128 px)}. The reach is the nominal lookup radius of the model:
its coarsest correlation level looks $32\,s$~px away from each point. A smaller $s$ resolves finer
detail, such as a strain concentration, and a larger $s$ reaches further, so pick the finest model
whose reach covers the largest displacement of any pair. In accumulative tracking that is the
largest displacement from frame~1; in incremental tracking it is the largest step between a frame
and its reference frame.

Hovering over an entry shows its description and reach, and \ui{Weights not installed.} when the
model file is not on the computer; the list checks this each time it opens, so a model you have
just downloaded or imported shows as installed. A model is downloaded the first time a run or a survey
uses it; \menu{Help}{Models…} shows which models are installed and imports model files on a
computer without an internet connection (Section~\ref{sec:installing-models}).

\subsection{Suggest a Model}
\label{sec:suggest}

\ui{Suggest a Model} measures how far the specimen moves and selects a model for you. Use it
before the first run, in particular when you do not know the size of the motion. It needs at
least two images and does not start while a run is going. It works in these steps:
\begin{enumerate}
  \item It loads RAFTcorr-s8, the model with the largest reach (downloading it on first use),
        with the current \ui{Device} and \ui{Iterations} (Section~\ref{sec:parameters}). A
        progress dialog, \ui{Measuring the motion with RAFTcorr-s8}, has a \ui{Cancel} button.
  \item It compares up to three pairs of frames, chosen by the tracking mode. In
        \ui{Accumulative} tracking these are frame~1 against the frame a quarter of the way
        through the sequence, the middle frame and the last frame (for 44 images, frames 11, 22
        and 44). In \ui{Incremental} tracking it takes the reference frames the run would use and,
        for each, the frame farthest from it that is compared with it, where the step is
        largest; of these pairs it surveys the last, the middle one and the one a quarter of the
        way through. The progress dialog names each pair, for example \ui{Survey: frame 44 against
        frame 43}.
  \item Over the region of interest of frame~1 (the whole image if none is drawn yet), it takes
        the 99th percentile of the displacement magnitude $\sqrt{u^2 + v^2}$ of each sampled
        pair. The 99th percentile ignores the largest 1\,\% of the values, so a few stray points
        do not decide the choice. The largest of the three values is the measured displacement
        $d$.
  \item It selects the finest model whose reach is at least $1.3\,d$: RAFTcorr-s1 for $d$ up to
        about 24.6~px, -s2 up to about 49.2~px, -s4 up to about 98.5~px, and -s8 above.
\end{enumerate}

The \ui{Model} list switches to the suggestion, and the same message appears in the \ui{LOG}
and in a card under the button (the parts in angle brackets are the measured values):
\begin{lstlisting}
Suggested: RAFTcorr-<s>. Largest displacement of the region of interest about <d> px (99th percentile over frames <a, b, c>); RAFTcorr-<s> reaches <R> px.
\end{lstlisting}
In incremental tracking the message lists the surveyed pairs, each as reference frame
$\rightarrow$ frame:
\begin{lstlisting}[literate={→}{{$\rightarrow$}}1]
Suggested: RAFTcorr-<s>. Largest displacement of the region of interest between a frame and its reference frame about <d> px (99th percentile over the pairs reference → frame <r1 → k1, r2 → k2, r3 → k3>); RAFTcorr-<s> reaches <R> px.
\end{lstlisting}
When $1.3\,d$ exceeds even the reach of RAFTcorr-s8 (above about 197~px), the suggestion is still
RAFTcorr-s8, and the message adds a sentence that starts with \ui{This exceeds even the reach of
RAFTcorr-s8:} and says what to change: \ui{use incremental tracking.} in accumulative tracking,
\ui{update the reference frame more often.} with \ui{Every N Frames} or \ui{Custom Frames}, and
\ui{the motion between consecutive frames is too large to track.} with \ui{Every Frame}. The card
stays until you open other images, and a project file keeps it.

\begin{tipbox}
Draw the region of interest before you click \ui{Suggest a Model}. The survey then measures the
motion of the specimen, not of the background around it.
\end{tipbox}

\begin{notebox}
The survey looks at three pairs only. If the motion peaks between the sampled frames and comes
back, or one step of an incremental run is much longer than the surveyed ones, the survey does
not see it. Choose the tracking mode before you click \ui{Suggest a Model}: a suggestion made
for accumulative tracking does not fit an incremental run, and the other way round.
\end{notebox}

\subsection{Which combination to choose}

\begin{center}
\begin{tabularx}{\textwidth}{>{\raggedright\arraybackslash}X
                             >{\raggedright\arraybackslash}p{4.2cm}
                             >{\raggedright\arraybackslash}p{3.6cm}}
  \toprule
  Situation & Tracking Mode & Model \\
  \midrule
  The displacement from frame~1 stays within the reach of a model in every frame
    & \ui{Accumulative} & the finest model that reaches it (\ui{Suggest a Model}) \\
  The displacement from frame~1 exceeds the reach of RAFTcorr-s8
    & \ui{Incremental}, \ui{Every Frame} & the finest model that reaches the step between
      two images \\
  Many images, each step small, large total motion
    & \ui{Incremental}, \ui{Every N Frames} & a model whose reach covers the motion over
      $N$ frames \\
  Uneven motion (slow, then fast)
    & \ui{Incremental}, \ui{Custom Frames} & a model whose reach covers the longest step \\
  \bottomrule
\end{tabularx}
\end{center}

\section{Region of interest}
\label{sec:roi}

The region of interest tells RAFTcorr which points to measure. It is a mask on frame~1, and
every result is reported on its points: the result arrays cover its bounding box, and the points
outside the region stay empty. A run does not start until frame~1 has one. The network is given
the bounding box of the region plus a margin (Section~\ref{sec:parameters}); a region drawn
tightly around the specimen keeps that box, and the memory each pair needs, small.

Any other frame can hold a mask of its own. Such a mask is drawn on the image of that frame,
where the specimen is at that moment, and it can leave out a crack or a hole there or serve as
the region of interest of a reference frame (Section~\ref{sec:roi-later-frames}).

\subsection{The hint}

The box at the top of the \ui{REGION OF INTEREST} section says which frames need a region of
interest. It follows the tracking mode (Section~\ref{sec:workflow}):
\begin{itemize}
  \item Before any images are loaded: \ui{Load images first, then draw a Region of Interest on
        frame 1.}
  \item \ui{Accumulative mode}: only frame~1 needs a region of interest. \ui{All later frames
        are compared against it directly.}
  \item \ui{Incremental}: only frame~1 needs a region of interest. RAFTcorr carries it to the
        reference frames, which the hint lists (the first eight, then an ellipsis), with the
        measured displacement. \ui{Draw on another frame only to replace that region or to
        leave out cracks and holes.}
\end{itemize}

\subsection{The toolbar}

\screenshot[0.45\textwidth]{fig-roi-toolbar.png}{The \ui{REGION OF INTEREST} section with the
hint, the toolbar and the menu of \ui{+ Add} open: \ui{Polygon}, \ui{Rectangle}, \ui{Circle} and
\ui{Circle (3-point)}.}{fig:roi-toolbar}

Below the hint are three rows of buttons (Figure~\ref{fig:roi-toolbar}): \ui{+ Add} and
\ui{Cut} (the button with the scissors); \ui{Import} and \ui{Batch Import}; \ui{Save},
\ui{Invert} and \ui{Clear}. Except \ui{Batch Import}, every button acts on the mask of the frame
shown on the canvas, and drawing, importing or clearing opens that frame for editing
(Section~\ref{sec:roi-editing}).

\subsubsection*{Drawing shapes}

\ui{+ Add} and \ui{Cut} open the same menu of four shapes. A shape picked from \ui{+ Add} is
added to the mask; a shape picked from \ui{Cut} is removed from it, which is how you leave out
holes, voids and damaged areas. Then draw the shape on the canvas:
\begin{description}[leftmargin=1.2cm,style=nextline]
  \item[Polygon]
    Click each vertex in turn; double-click to place the last vertex and close the polygon, or
    press \kbd{Enter} to close it after the last vertex you clicked.
    A polygon needs at least three vertices; with fewer, nothing is drawn.
  \item[Rectangle]
    Press at one corner, drag to the opposite corner and release.
  \item[Circle]
    Press at the centre, drag outward to set the radius and release.
  \item[Circle (3-point)]
    Click three points on the edge of the circle. After the second click a preview circle runs
    through the two points and the cursor; the third click fixes the circle that passes through
    all three. Use it when the edge of a round feature is visible but its centre is not. If the
    three points lie almost on one line, no circle is drawn and the log asks for points spread
    around the edge.
\end{description}

While you draw, the cursor is a cross, the outline is blue for \ui{+ Add} and red for \ui{Cut},
and the button you started from is highlighted. A shape is applied as soon as it is complete.
Each pick from the menu draws one shape; pick again for the next one. If you pick the wrong
shape, pick the right one; it replaces the first. To abandon a shape, press \kbd{Esc} before it
is complete, whether you have started drawing it or only picked it from the menu. Shapes that reach beyond the image are cut at its border.

\begin{notebox}
The toolbar has no brush and no undo. Before a change you may want to take back, click
\ui{Save} to keep a copy of the mask, and \ui{Import} it to return to it.
\end{notebox}

\subsubsection*{Import, Save, Invert and Clear}

\begin{description}[leftmargin=1.2cm,style=nextline]
  \item[Import]
    Opens the \ui{Import Mask Image} dialog (PNG, BMP, TIFF or JPEG). The file replaces the mask
    of the frame shown. Bright pixels are inside: a pixel counts as inside when it is brighter
    than mid-grey (127 on a scale of 0 to 255), so masks stored as 0 and 255 and masks stored as
    0 and 1 both work. A colour file is read as its luminance. A mask of another size than the
    images is resized to the image size (nearest neighbour). If the file cannot be read, the
    mask stays as it was, and the log says \emph{``Could not read the mask file''} with the file
    name and the reason.
  \item[Batch Import]
    Assigns mask files to many frames at once (Section~\ref{sec:roi-batch}).
  \item[Save]
    Writes the mask of the frame shown to a PNG file, white (255) inside and black (0) outside.
    The dialog proposes a name made of the image folder and the frame number, next to the image
    folder rather than inside it: for frame~1 of the folder \file{specimen}, the file
    \file{specimen\_roi\_frame1.png} beside \file{specimen}. A mask saved among the images would be
    loaded as an image the next time you open the folder. \ui{Import}, \ui{Batch Import}
    and \file{raftcorr run --roi} (Section~\ref{sec:cli}) read this format.
  \item[Invert]
    Swaps inside and outside of the mask of the frame shown, whether you drew it or it came from
    \ui{Import}, \ui{Batch Import}, the image list or a project file. On frame~1 without a mask it
    selects the whole image. On a later frame without a mask of its own, while frame~1 has one,
    nothing changes, and the log explains that the region of frame~1 is used there.
  \item[Clear]
    Removes the mask of the frame shown. The masks of the other frames stay; to clear several
    frames at once, select them in the image list and right-click (Section~\ref{sec:images}).
\end{description}

\subsection{Editing mode}
\label{sec:roi-editing}

\screenshot[0.5\textwidth]{fig-roi-canvas.png}{The canvas in editing mode for frame~1: the
banner \ui{EDITING REGION OF INTEREST FOR FRAME 01} above the image, and the region of interest
in blue over the gauge section of a plate with two holes, drawn as a rectangle with two circles
cut out.}{fig:roi-canvas}

While you edit a mask, a yellow banner (amber in the light theme) above the canvas names the
frame: \ui{EDITING REGION OF INTEREST FOR FRAME 01} (Figure~\ref{fig:roi-canvas}). The canvas
shows the image of that frame with its mask as a blue tint.

You enter editing mode with the \ui{Region} button of a frame in the image list (\ui{Need},
\ui{Edit} or \ui{Add}), which also shows that frame, or by drawing, importing or clearing with
the toolbar, which edits the frame shown. To edit another frame, click its \ui{Region} button, or move to it with
the image list or the frame slider: the banner follows. Every change is stored at once; there is
nothing to confirm. Editing mode ends when a run starts. Once results exist, moving to another
frame also ends it, and the canvas shows the displacement field again.

A typical region of interest for a plate with two holes (Figure~\ref{fig:roi-canvas}):
\begin{enumerate}
  \item Choose the tracking mode first (Section~\ref{sec:workflow}).
  \item Show frame~1 and choose \menu{+ Add}{Rectangle}. Drag a rectangle over the gauge
        section.
  \item Choose \menu{Cut}{Circle} and drag a circle over each hole, or use
        \menu{Cut}{Circle (3-point)} and click three points on the edge of each hole.
  \item Check the blue region. Click \ui{Save} to keep a copy for later runs or for the command
        line.
\end{enumerate}

\subsection{Masks on later frames}
\label{sec:roi-later-frames}

A mask on a frame other than frame~1 is drawn, imported or batch-imported on the image of that
frame, in its own coordinates. RAFTcorr reads the mask of every frame that has one
(Table~\ref{tab:frame-masks}). In every tracking mode, a point of the region of interest whose
position in frame $k$ falls outside the mask of frame $k$ is left empty in frame $k$: this is the
way to remove a crack, a hole that opens, or a part of the image where the speckle is lost. In
incremental tracking, the mask of a reference frame also replaces the region of frame~1 carried
to that frame, and the pairs compared with it are solved over the mask.

\begin{table}[htbp]
  \centering
  \caption{What the mask of frame $k$ (not frame~1) does.}
  \label{tab:frame-masks}
  \begin{tabularx}{\textwidth}{>{\raggedright\arraybackslash}p{3.2cm}
                               >{\raggedright\arraybackslash}X
                               >{\raggedright\arraybackslash}X}
    \toprule
    Tracking Mode & Frame $k$ has a mask & Frame $k$ has no mask \\
    \midrule
    \ui{Accumulative}
      & Points that fall outside it are empty in frame $k$ only.
      & Nothing is left out. \\
    \ui{Incremental}, frame $k$ is a reference frame
      & Points outside it are empty in frame $k$ and in every frame measured from it; the pairs
        compared with frame $k$ are solved over the mask.
      & The region of frame~1 is carried to frame $k$ with the displacement measured so far. \\
    \ui{Incremental}, frame $k$ is not a reference frame
      & Points that fall outside it are empty in frame $k$.
      & Nothing is left out. \\
    \bottomrule
  \end{tabularx}
\end{table}

Three rules follow from this:
\begin{itemize}
  \item \textbf{A mask on a later frame covers everything that stays.} It is not a mask of the
        crack; it is the whole specimen as it appears in that frame, minus the crack. On a frame
        without a mask of its own, drawing starts from an empty mask, so a \ui{Cut} alone leaves
        nothing. Add a shape over the specimen first, then cut the crack out, or import masks
        made by a segmentation tool with \ui{Batch Import}.
  \item \textbf{A mask cannot add points.} Results exist only for the points of the region of
        interest of frame~1. A mask on a later frame that is larger than the specimen adds
        nothing.
  \item \textbf{In accumulative tracking, a mask acts on its own frame only.} A crack that stays
        open needs a mask on every frame where it is open. \ui{Batch Import} assigns one mask
        file to many frames at once.
\end{itemize}

\begin{tipbox}
In incremental tracking the region of frame~1 follows the specimen from one reference frame to
the next. A specimen that moves or rotates across the image does not need masks on the
reference frames; give one only where you want to leave something out.
\end{tipbox}

\subsection{Batch Import}
\label{sec:roi-batch}

\screenshot[\textwidth]{fig-batch-import.png}{The \ui{Batch Import Region of Interest Masks}
dialog: the \ui{Available Masks} of the chosen folder (left), the \ui{Frame Assignments} with
their \ui{Frame}, \ui{Image} and \ui{Mask} columns (middle), and the \ui{Preview} of the selected
frame with its mask (right).}{fig:batch-import}

\ui{Batch Import} opens the \ui{Batch Import Region of Interest Masks} dialog
(Figure~\ref{fig:batch-import}). It lists every frame, because RAFTcorr reads the mask of every
frame.
\begin{description}[leftmargin=1.2cm,style=nextline]
  \item[Mask Folder]
    \ui{Browse...} chooses the folder of the mask files. Its image files (PNG, BMP, TIFF, JPEG,
    JPEG~2000 and WebP) appear under \ui{Available Masks}, in natural order (\file{mask2.png}
    before \file{mask10.png}). A mask whose size differs from the
    images is greyed out and cannot be assigned; its tooltip gives both sizes, and a line above
    the lists counts these masks. Unlike \ui{Import}, the dialog does not resize masks.
  \item[Auto-Match by Name]
    Gives each image the mask whose file name ends in the same number: \file{mask6.png} goes to
    \file{image6.tif}, whatever frame that image is. If several images end in that number, each
    of them gets the mask. Only when the image names carry no number is the last number of the
    mask name read as the frame position counted from 0: \file{mask0.png} goes to frame~1.
  \item[Assign Sequential]
    Gives the masks, in the order of the list, to frames 1, 2, 3, and so on.
  \item[Assign Selected ->]
    Pairs the masks selected on the left with the frames selected in the middle: one mask with
    one frame, one mask with many frames, or as many masks as frames in the order of selection.
    Several masks for one frame are refused.
  \item[Clear All]
    Removes every assignment. \ui{Auto-Match by Name} and \ui{Assign Sequential} also start
    from no assignments; \ui{Assign Selected ->} adds to the existing ones.
  \item[Preview]
    Shows the image of the frame selected in the middle with its assigned mask, or the mask
    selected on the left over that image. \ui{View:} switches between \ui{Image only},
    \ui{Image + Mask} (default) and \ui{Mask only}; \ui{Alpha:} sets the opacity of the mask
    on a slider from 0 to 100 (default 50); \ui{Mask color:} offers \ui{Blue}, \ui{Red},
    \ui{Green} and \ui{Yellow}. The mouse wheel zooms and dragging pans.
\end{description}
\ui{OK} reads the assigned files with the same rule as \ui{Import} (bright is inside) and
replaces the masks of those frames; the log lists every frame that received one. \ui{Cancel}
changes nothing.

\begin{warnbox}
When the image names carry no number, \ui{Auto-Match by Name} counts frames from 0, as data
files do, and masks numbered from 1 land one frame late. Use \ui{Assign Sequential} or
\ui{Assign Selected ->} instead, and check the \ui{Mask} column before you click \ui{OK}.
\end{warnbox}

A mask with no white pixel, assigned to frame~1, leaves frame~1 without a region of interest:
the run then does not start (Section~\ref{sec:running}).

\section{Parameters}
\label{sec:parameters}

Two sections below \ui{REGION OF INTEREST} hold the remaining settings of a run
(Figure~\ref{fig:parameters}). \ui{PARAMETERS},
open by default, sets how the network runs; \ui{ADVANCED}, closed by default (click its title to
open it), sets how each image pair is cut into windows for the network. Every new session starts
from the defaults in Table~\ref{tab:parameters}; a project file keeps the values you set
(Section~\ref{sec:projects}).
% CHECK: after File > Open Project the ADVANCED controls and the Interval box keep their previous values on screen, although the run uses the project's values (AdvancedPanel has no sync from AppState); reproduced offscreen

\begin{table}[htbp]
  \centering
  \caption{The controls of \ui{PARAMETERS} and \ui{ADVANCED}.}
  \label{tab:parameters}
  \begin{tabularx}{\textwidth}{>{\raggedright\arraybackslash}p{2.6cm}
                               >{\raggedright\arraybackslash}p{5.2cm}
                               >{\raggedright\arraybackslash}X}
    \toprule
    Control & Default & Range or choices \\
    \midrule
    \ui{Iterations}   & 20 iterations & 1 to 100 iterations \\
    \ui{Smoothing}    & \ui{Off} (0 px) & 0 to 50 px, in steps of 0.5 px \\
    \ui{Device}       & \ui{Automatic} & \ui{Automatic}, \ui{GPU (CUDA)}, \ui{CPU} \\
    \midrule
    \ui{Margin}       & \ui{Use the Model Reach} ticked ($32\,s$ px) & 0 to 4096 px \\
    \ui{GPU Memory}   & 0.85 (fraction of the free memory) & 0.05 to 1.00, in steps of 0.05 \\
    \ui{Tile Blending} & 32 px & 0 to 512 px \\
    \ui{Correlation}  & \ui{Automatic} & \ui{Automatic}, \ui{On Demand}, \ui{Stored Volume} \\
    \bottomrule
  \end{tabularx}
\end{table}

\screenshot[0.45\textwidth]{fig-parameters.png}{The \ui{PARAMETERS} and \ui{ADVANCED} sections,
both open, with their default values.}{fig:parameters}

Hovering over a label shows what it does. A spin box or a list reacts to the mouse wheel only
after you click it, so scrolling the sidebar does not change a value by accident.

\subsection{PARAMETERS}

\begin{description}[leftmargin=1.2cm,style=nextline]
  \item[Iterations]
    The number of refinement iterations of the network (default 20). The network refines its
    estimate of the displacement once per iteration, so the computing time grows with the
    number. \ui{Suggest a Model} uses the same value, and the card on the canvas shows it under
    \ui{ITERATIONS}. Keep the default unless you have a reason to change it; if you lower it to
    shorten a run, compare the residual (Section~\ref{sec:results}) with that of a run at the
    default before you rely on the result.
  \item[Smoothing]
    Gaussian smoothing of the tracked displacement, set as its standard deviation $\sigma$ in
    pixels. It is \ui{Off} by default, so the results hold the displacement as the network
    measures it. Smoothing coarsens the spatial resolution: features of the field smaller than
    a few $\sigma$ are flattened. The smoothing stays inside the region of interest: empty points
    and points outside the region are not used and stay empty, and a point with almost no valid
    neighbours (less than 5\,\% of the weight of the Gaussian) becomes empty. The smoothed
    displacement is what the results, the strain window and the exports receive; in
    incremental tracking it is also what later frames are composed with. For strain alone,
    prefer \ui{Pre-fit smoothing} in the strain window (Section~\ref{sec:strain}), which feeds
    only the strain fit and leaves the displacement as it is.
  \item[Device]
    Where the network runs. \ui{Automatic} uses the GPU when PyTorch sees one (an NVIDIA GPU
    with CUDA) and the CPU otherwise; on macOS RAFTcorr runs on the CPU. \ui{GPU (CUDA)} on a
    computer without CUDA stops the run with an error. Choose \ui{CPU} to leave the GPU to
    another program. \menu{Help}{Self-test} checks the installed models on the device you choose
    (Section~\ref{sec:installing-check}).
\end{description}

\subsection{ADVANCED}

For every pair, the network receives a \emph{window} of both images: the bounding box of the
region of interest of the reference frame, grown by a margin on every side. When that window
needs more memory than the device offers, RAFTcorr cuts the region into \emph{tiles}. The plan is
automatic; Section~\ref{sec:tiling} describes it, its line in the \ui{LOG} and what happens when the
memory runs out. The four controls below adjust it.

\begin{description}[leftmargin=1.2cm,style=nextline]
  \item[Margin and Use the Model Reach]
    The margin is the image context added around the bounding box, so that a point that moves
    stays inside the window the network sees. It must cover the largest displacement of the
    region of interest. With \ui{Use the Model Reach} ticked (default), the margin is the reach
    of the selected model: 32, 64, 128 or 256~px for RAFTcorr-s1, -s2, -s4 or -s8; the box is
    greyed out and its value is not used. Untick it to type a margin from 0 to 4096~px. A
    smaller margin than the reach saves memory when you use a model with a large $s$ on small
    motion; a margin smaller than the displacement loses the points that move out of the
    window. The window never extends past the image, and when tiles are needed the margin can
    be reduced to fit the memory, in which case the tiling line of the log says so.
  \item[GPU Memory]
    The fraction of the free GPU memory one window may use, measured before each pair (default
    0.85; RAFTcorr also keeps 384~MB free). A larger region is cut into tiles to fit. Lower it
    when another program needs the GPU or when runs fail for lack of memory; raise it when the
    GPU is free and you want fewer tiles. On the CPU it has no effect: windows are planned for a
    fixed budget of 6~GB.
  \item[Tile Blending]
    The width of the overlap between neighbouring tiles, where their results are blended with
    linear weights (default 32~px). It matters only when the region is cut into tiles. A wider
    overlap spreads the change from one tile to the next over more pixels and needs more tiles;
    0~px lets the tiles meet edge to edge. The default rarely needs a change.
  \item[Correlation]
    How the network computes its correlation lookup. Both ways give the same displacement; they
    differ in memory and speed. \ui{Stored Volume} computes the correlation between all pairs of
    feature points of the window once and keeps it, so its memory grows with the square of the
    window area; it is the faster choice where it fits. \ui{On Demand} computes each lookup when it is needed,
    with memory in proportion to the window area. \ui{Automatic} uses \ui{Stored Volume} for
    RAFTcorr-s8 and \ui{On Demand} for -s1, -s2 and -s4. Leave it at \ui{Automatic}; with
    RAFTcorr-s8 and little memory, \ui{On Demand} allows larger windows and fewer tiles.
\end{description}

\section{The canvas}
\label{sec:canvas}

The canvas fills the centre of the main window (callout~3 in Figure~\ref{fig:main-window}). It
shows the current frame: before a run, the image with its mask as a blue tint; while you edit a
mask, the same with the editing banner (Section~\ref{sec:roi-editing}); after a run, the
displacement field over the image, with a colorbar (Section~\ref{sec:results}). A toolbar runs along its
top edge and the frame navigator along its bottom edge.

\subsection{Zoom and pan}

\begin{description}[leftmargin=1.2cm,style=nextline]
  \item[Fit]
    Fits the whole image into the canvas. RAFTcorr also fits the image when you open a folder.
  \item[100\%]
    Shows one image pixel per screen pixel.
  \item[Zoom in, Zoom out]
    The two magnifier buttons next to \ui{100\%} zoom by one step, a factor of 1.15.
  \item[Mouse wheel]
    Zooms by the same step around the point under the cursor, which stays in place.
  \item[Middle mouse button]
    Drag with the middle button (the pressed wheel) to pan. The canvas has no scroll bars.
\end{description}
The wheel and the zoom buttons stop at 10\,\% and 2000\,\%. Without a middle button, as on a touchpad, move
around by zooming out (or \ui{Fit}) and zooming in again over the place you want to see.

At the right end of the toolbar, a read-out follows the cursor: its position as
\file{x 512.0  y 384.0 px}, and, when a field is shown, the value of the field under the cursor
with its unit. The position is in image coordinates, $x$ to the right and $y$ downward in pixels,
with the centre of the top-left pixel at $(0, 0)$, the same convention as the data files.

\subsection{The configuration card}

Once at least two images are loaded, a card in the top-left corner of the canvas summarises what
a run would do:
\begin{description}[leftmargin=1.2cm,style=nextline]
  \item[MODE] \ui{Accumulative}, \ui{Incremental}, \ui{Incremental, every N frames} (with the
    interval) or \ui{Custom reference frames} (Section~\ref{sec:workflow}).
  \item[MODEL] The selected model, for example \ui{RAFTcorr-s4}.
  \item[ITERATIONS] The number of iterations (Section~\ref{sec:parameters}).
\end{description}
The card follows the sidebar as you change the settings. Mouse clicks pass through it to the
canvas.

\subsection{The frame navigator}

The bar below the canvas holds, from left to right:
\begin{description}[leftmargin=1.2cm,style=nextline]
  \item[Previous frame, Play animation, Next frame]
    The arrows step one frame back or forward and stop at the first and the last frame. Play
    advances one frame at a time and starts again at frame~1 after the last frame; the same
    button pauses. Opening other images stops the playback.
  \item[Playback speed]
    \file{1 fps}, \file{2 fps} (default), \file{5 fps}, \file{10 fps} or \file{30 fps}. This is
    the speed of the display, not the frame rate of the experiment; that one belongs to the
    physical units (Section~\ref{sec:physical-units}).
  \item[FRAME n/N]
    The current frame and the number of frames, counted from 1 as in the image list.
  \item[Slider]
    Drag it to move through the sequence.
\end{description}
The slider, the image list and the canvas always show the same frame: moving one moves the
others.

RAFTcorr has no keyboard shortcut for frames that works everywhere in the window. \kbd{Up},
\kbd{Down}, \kbd{Page Up}, \kbd{Page Down}, \kbd{Home} and \kbd{End} move through the frames when
the image list has the keyboard focus (click a row first), and the same keys move the slider after
you click it. Table~\ref{tab:shortcuts} lists every key the application answers to.
% CHECK: slider keys are Qt's standard QSlider behaviour (singleStep 1, pageStep 10 by default), not RAFTcorr code; confirm on Windows and macOS

\section{Running the analysis}
\label{sec:running}

The top of the right sidebar holds the four buttons that start, stop and use a run, and the
\ui{PROGRESS} area below them (Figure~\ref{fig:progress}). This section follows a run from the
click on \ui{Run DIC Analysis} to the moment its results appear.

\screenshot[0.42\textwidth]{fig-progress.png}{The top of the right sidebar during a run:
\ui{Run DIC Analysis} (disabled while the run goes on), \ui{Cancel}, \ui{Export Results},
\ui{Open Strain Window} and the \ui{PROGRESS} area with its bar, the status line and the
\ui{ELAPSED} and \ui{REMAINING} times.}{fig:progress}

\subsection{Run DIC Analysis}

\ui{Run DIC Analysis} tracks every frame of the sequence with the tracking mode, the model and
the parameters set in the left sidebar (Sections~\ref{sec:workflow} and~\ref{sec:parameters}).
The button is enabled whenever no run is going on. Before anything is computed, RAFTcorr checks
the inputs; a missing input opens a \ui{Cannot start analysis} dialog that says what to do:

\begin{itemize}
  \item fewer than 2 images: \emph{Need at least 2 images. Load an image folder in the IMAGES
        section before running DIC.}
  \item no region of interest on frame~1: \emph{Define a Region of Interest first: draw or
        import one for frame 1 in the REGION OF INTEREST section.}
  \item a region of interest on frame~1 with no pixel selected, such as an all-black mask from
        \ui{Batch Import}: \emph{The Region of Interest of frame 1 is empty (no pixels selected).
        Draw it again in the REGION OF INTEREST section.}
  \item refused text in the \ui{Reference Frames} field, with \ui{Custom Frames}: the dialog
        repeats the reason the field gives (Section~\ref{sec:workflow});
  \item a setting the run cannot use: the same dialog names it.
\end{itemize}

In incremental tracking, a reference frame without a region of interest of its own does not stop
the run: RAFTcorr carries the region of frame~1 to it, and the \ui{LOG} lists such frames
(Section~\ref{sec:workflow}).

The run then starts in the background. The \ui{LOG} records what it does: the model, the number
of frames, the tracking mode, the number of iterations and the device; the line \emph{loaded
RAFTcorr-s4 (\ldots)} when the model is loaded into memory (it stays loaded for the next run with
the same model and device); the size and position of the region of interest; how the region is
cut into windows (Section~\ref{sec:tiling}). A frame whose tracked part covers less
than 90\,\% of the region of interest gets a warning line, \emph{Frame $k$: only \ldots{} of the
region of interest could be tracked}.

Results appear only when the run ends: there is no preview of frames during the run. Until then
the canvas keeps showing the previous results, if any. They are replaced only when the new run
returns at least one computed frame. Results larger than 2\,GB are kept on disk in a temporary
folder until you quit (\emph{results of \ldots{} GB are kept on disk in \ldots},
Section~\ref{sec:overview-files}).

Loading other images while a run is going on cancels the run (\emph{The images changed: the
running analysis is cancelled.}), because its results would not belong to the new images.

\subsection{Cancel}

\ui{Cancel} is enabled while a run is going on. It stops the run at the next check point:
between two frames, between two tiles of a frame, or during a model download. The run does not
stop in the middle of a window, so it can take a moment.

Frames already computed are kept. The run returns them as its results, as a complete run would:
the \ui{LOG} reports \emph{stopped by the user after $k$ of $n$ pairs} and then \emph{Results
received: $k$ frame(s)}; \ui{Export Results} is enabled, the strain window opens, and the
tracking-quality check (Section~\ref{sec:tracking-quality}) runs on the computed frames. The
frames that were not computed have no result: the canvas shows their image without a field, the
strain window shows \emph{Frame $k$ was not computed.}, and exported data files hold NaN for them
(with a \file{Done} array that marks the computed frames, Section~\ref{sec:export}).

If you cancel before the first frame is finished, nothing replaces the previous results
(\emph{No frame was computed; the previous results are kept.}), and the strain window does not
open. Cancelling during the model
download stops the download; the \ui{LOG} says \emph{Stopped.}

\begin{warnbox}
A cancelled run that finished at least one frame replaces the previous results with its partial
ones. Save a project (\menu{File}{Save Project}, \kbd{Ctrl+S}) before you start a run you may
cancel if you still need the previous results.
\end{warnbox}

A run that stops on an error is different: it returns no frames at all. An \ui{Analysis failed}
dialog shows the cause, the \ui{LOG} holds the full traceback, the previous results, if any,
stay, and the strain window does not open.

\subsection{The PROGRESS area}

\begin{description}[leftmargin=1.2cm,style=nextline]
  \item[Progress bar]
    The fraction of the frame pairs done. It also moves inside a pair when the pair is cut into
    tiles.
  \item[Status line]
    \ui{Ready} when no run is going on. During a run it shows the percentage done and
    \ui{Frame $k$/$N$}: $k$ is the last frame finished and $N$ the number of images, frame~1
    included. Both count from 1, like the frame slider. Frame~1 is the reference, so the first
    frame to finish is frame~2 or later. Before the first frame finishes, the line shows the
    percentage alone.
  \item[Model download]
    While a run downloads its model (first use only), the status line shows \emph{Downloading
    model RAFTcorr-s4: $a$ of $b$ MB} instead of the frame, and the bar follows the download. A
    model that can be neither found nor downloaded stops the run and opens the Models window
    (Section~\ref{sec:installing-models}).
  \item[ELAPSED]
    The time since you clicked \ui{Run DIC Analysis}, in minutes and seconds, updated every
    second. It includes loading, and downloading, the model.
  \item[REMAINING]
    An estimate from the elapsed time and the fraction done; \ui{-{}-:-{}-} until 1\,\% of the
    run is done. Because the elapsed time includes loading the model, the first estimates are
    often too long; the estimate settles after a few frames.
\end{description}

\subsection{Large regions: tiles}
\label{sec:tiling}

The network sees a window of the reference image and the same window of the deformed image. A
point can only be matched if its displaced position lies inside that window, so the window is
the bounding box of the region of interest grown by a margin. By default the margin is the reach
of the model, 32\,$s$\,px (128\,px for RAFTcorr-s4); \ui{Margin} and \ui{Use the Model Reach}
under \ui{ADVANCED} change it (Section~\ref{sec:parameters}).

The memory one pass of the network needs grows with the number of pixels in the window, at a
rate that depends on the model: a larger $s$ needs less memory per pixel, and the stored
correlation volume that RAFTcorr-s8 uses by default grows faster than the window area. Before
each frame pair, RAFTcorr estimates the peak memory of one pass with a memory model of the
selected model and compares it with the share of the free GPU memory set by \ui{GPU Memory}
(default 0.85, under \ui{ADVANCED}), or with a fixed 6\,GB on the CPU.

\begin{itemize}
  \item If the whole window fits, the pair is solved in one window.
  \item If it does not, the region is cut into tiles. Each tile has a \emph{core}, the pixels it
        writes, and a \emph{halo} around it, the margin, which the network sees but whose result
        is not used. Neighbouring cores overlap by \ui{Tile Blending} (default 32\,px), and the
        overlap is blended with linear weights so that no seam shows.
\end{itemize}

Tiling is automatic and needs no setting. The \ui{LOG} says which case applies, at the first pair
and again whenever the plan changes (in incremental tracking the region of each reference frame
differs, and the free memory can change during a run):

\begin{itemize}
  \item \emph{tiling: one window of $W$ x $H$ px (margin $M$ px)};
  \item \emph{tiling: $n$ tiles, windows up to $W$ x $H$ px, margin $M$ px}.
\end{itemize}

\begin{warnbox}
When the memory only allows small windows, RAFTcorr reduces the halo so that each tile keeps a
useful core, and the tiling line ends with \emph{(margin reduced to fit memory)}. A point that
moves farther than the reduced margin can then leave the window of its tile and be matched
wrongly. Free GPU memory (close other GPU programs), raise \ui{GPU Memory}, choose a model with a
larger $s$, or reduce the region of interest.
\end{warnbox}

The memory model is an estimate. If the GPU still runs out of memory, RAFTcorr empties its memory
cache, makes the windows smaller (at most 55\,\% of the area of the largest window that failed) and
solves the pair again; the \ui{LOG} shows \emph{out of device memory; retrying with windows of at
most \ldots{} px}. After three such retries the run stops with an \ui{Analysis failed} dialog
that begins with \emph{The GPU ran out of memory. Close other GPU programs, choose a model with a
larger s, or lower GPU Memory in ADVANCED.} A budget so small that a window would fall below the
smallest size the model accepts also stops the run, with a message that gives both sizes and
begins with \emph{There is not enough free GPU memory for the smallest window this model accepts.
Close other GPU programs, or raise GPU Memory in ADVANCED if it is set low.}

\subsection{Tracking quality warning}
\label{sec:tracking-quality}

When a run ends, RAFTcorr checks the residual of every computed frame: the grey-level mismatch
that remains after the deformed image is mapped back onto frame~1 with the measured displacement
(Section~\ref{sec:residual}). It takes the median residual over the region of interest of each
frame. The residual grows slowly as the specimen deforms, but it jumps when tracking fails,
typically when the displacement from the reference frame exceeds the reach of the model.

A frame counts as a jump when its median residual exceeds 1.6 times that of the previous computed
frame and 2 times that of the first computed frame. From the first jump on, every frame whose
median residual stays above 1.6 times the level before the jump is flagged. The check needs at
least two computed frames after frame~1, and it runs after a cancelled run too, on the frames
computed.

When frames are flagged, a \ui{Tracking Quality Warning} dialog opens with the summary \emph{Part
of the sequence may not be tracked correctly.}, then the finding and an advice that depends on the
tracking mode:

\begin{itemize}
  \item \emph{The tracking quality drops from frame $k$ on: the intensity residual rises to $r$
        grey levels, against $r_0$ before.} Here $r_0$ is the highest median residual of the
        frames before frame~$k$.
  \item In accumulative tracking: \emph{The displacement from frame 1 may exceed the reach of the
        model. Try incremental tracking, or a model with a larger s.}
  \item In the other tracking modes: \emph{The step between reference frames may exceed the reach
        of the model. Use more reference frames, or a model with a larger s.}
\end{itemize}

\ui{Show Details\ldots} lists the flagged frames (the first ten, then how many more) and suggests
showing the \ui{Residual} field. Frame numbers count from 1, as on screen. The dialog does not
block the application: it sits in front of the strain window, you can keep working while it is
open, and the next run closes it. The \ui{LOG} keeps the same finding in one line
(\emph{tracking quality drops from frame $k$ on \ldots}), and the note is stored with the
provenance of the run in exported files.

The check never changes a result. When it fires:

\begin{enumerate}
  \item Select \ui{Residual} in \ui{FIELD} and step through the flagged frames.
  \item If the residual is high over most of the region, tracking was lost. In accumulative
        tracking, switch \ui{Tracking Mode} to \ui{Incremental}, or choose a model with a larger
        $s$ (\ui{Suggest a Model} measures the motion and picks one,
        Section~\ref{sec:workflow}). In incremental tracking, use more reference frames
        (\ui{Every Frame}, a smaller \ui{Interval}, or more \ui{Reference Frames}), or a model
        with a larger $s$.
  \item If the residual is high only where the images themselves change (a crack that opens,
        light reflected from the surface, speckle that flakes off), the displacement elsewhere
        can still be right. The median used by the check moves only when such an area is large;
        exclude it with a region of interest on the later frames (Section~\ref{sec:roi}).
\end{enumerate}

\subsection{When a run finishes}

At the end of a run that returns at least one computed frame:

\begin{itemize}
  \item the \ui{LOG} shows \emph{done: $n$ frames in \ldots{} s} and, in green, \emph{Results
        received: $n$ frame(s)}, where $n$ counts the computed frames, frame~1 excluded;
  \item \ui{Run DIC Analysis} is enabled again, \ui{Cancel} is disabled and \ui{Export Results}
        becomes available (also as \menu{File}{Export Results\ldots}, \kbd{Ctrl+E});
  \item the canvas moves to the last computed frame and draws the field selected in \ui{FIELD}
        (Section~\ref{sec:results});
  \item the strain window opens by itself, in front of the main window
        (Section~\ref{sec:strain}). It computes nothing until you click \ui{Compute Strain}. If
        you only need displacement, close it; \ui{Open Strain Window} brings it back;
  \item a strain result of the previous run is discarded, since it belongs to other
        displacements;
  \item the \ui{Tracking Quality Warning} dialog opens if frames were flagged.
\end{itemize}

\section{Viewing displacement results}
\label{sec:results}

After a run, the canvas draws the selected field over the image, with a colorbar at its right
edge (Figure~\ref{fig:displacement}). Three groups of the right sidebar decide what is drawn and
how: \ui{FIELD}, \ui{VISUALIZATION} and \ui{PHYSICAL UNITS} (Figure~\ref{fig:display}). The
\ui{LOG} at the bottom of the sidebar records what the application does.

\screenshot[0.42\textwidth]{fig-display.png}{The \ui{FIELD}, \ui{VISUALIZATION} and
\ui{PHYSICAL UNITS} groups of the right sidebar after a run.}{fig:display}

Every field is Lagrangian: a value belongs to a material point and is the change from frame~1 to
the frame on show. Frame~1 is the reference, so its displacement and residual are zero.

\subsection{FIELD}

Four buttons choose the field; one is always selected.

\begin{description}[leftmargin=1.2cm,style=nextline]
  \item[Disp U]
    Horizontal displacement $u$, positive to the right.
  \item[Disp V]
    Vertical displacement $v$, positive downward, as image rows count.
  \item[Magnitude]
    $\sqrt{u^2 + v^2}$, the length of the displacement.
  \item[Residual]
    The grey-level mismatch left after the deformed image is mapped back onto frame~1 with the
    measured displacement: a map of how well each point was tracked.
\end{description}

Displacements are in pixels, or in a physical length unit when \ui{Use physical units} is on
(Section~\ref{sec:physical-units}). The residual is in grey levels.

\screenshot[0.45\textwidth]{fig-displacement.png}{Vertical displacement \ui{Disp V} of a plate
with two holes in tension, drawn on the deformed last frame (\ui{Show on}: \ui{Deformed frame})
with the background image. The colorbar gives the field and its unit, V~(px), and the range of
the colours.}{fig:displacement}

\subsubsection*{The residual}
\label{sec:residual}

For each point $\mathbf{X}$ of the region of interest, RAFTcorr reads the current frame at the
displaced position $\mathbf{X} + \mathbf{u}(\mathbf{X})$ and subtracts the frame-1 image at
$\mathbf{X}$. It removes the local mean of this difference and takes its local root mean square,
both over a Gaussian window of $\sigma = 4$\,px. The result is in grey levels on the 0 to 255
scale to which RAFTcorr reads every image (Section~\ref{sec:images}). It is always computed
against frame~1 with the total displacement, whatever the tracking mode.

\begin{itemize}
  \item Where the displacement maps the current frame back onto frame~1, the residual is small,
        at the level of the image noise.
  \item It grows slowly as the specimen deforms (the speckles change shape) and with changes of
        lighting. Because the local mean is removed, a change of brightness that varies slowly
        across the image does not raise it.
  \item It is large where tracking failed, and where the images themselves differ: a crack that
        opens, a reflection, speckle that flakes off, a surface that leaves the focal plane.
\end{itemize}

Read the residual by comparison, between frames and between areas of one frame, rather than
against a fixed value. A residual that rises over the whole region from one frame to the next
means tracking was lost there; this is what the tracking-quality check detects
(Section~\ref{sec:tracking-quality}). A residual that rises in one area points to that area.
Points without a displacement have no residual. The default colormap of the residual is
\file{magma} (Figure~\ref{fig:residual}).

\screenshot[0.45\textwidth]{fig-residual.png}{The residual of the frame of
Figure~\ref{fig:displacement}, in grey levels.}{fig:residual}

\subsection{Where the field is drawn}

\begin{description}[leftmargin=1.2cm,style=nextline]
  \item[Show on]
    \ui{Deformed frame} (the default) draws each value where its material point is in the
    current frame, over the image of that frame. To do so, RAFTcorr inverts the measured
    displacement: for every pixel of the current frame it finds the point of frame~1 that lands
    there and draws that point's value. Pixels where no point of the region of interest lands
    stay empty. \ui{Reference frame} draws each value at the frame-1 position of its point,
    over the image of frame~1. The values are the same; only their positions differ. On
    frame~1 both views are the same.

    Use \ui{Deformed frame} to check that the field follows the specimen: its edges and holes
    should match the image underneath. Use \ui{Reference frame} to compare frames point by point:
    a point stays at one place on the screen while you play the sequence.
  \item[Show background image]
    On by default. Untick it to see the field alone. \ui{Show on} still applies: the positions
    of the points differ between the two views even with nothing behind them.
  \item[Hidden background]
    What replaces the image while it is hidden: \ui{White} (the default), \ui{Black} or
    \ui{Transparent}, which shows the background of the canvas. The list is enabled only while
    \ui{Show background image} is off. In exported pictures, transparency is kept for PNG and
    TIFF; other formats get white.
  \item[Hide untracked regions]
    On by default. A point of the region of interest is untracked in a frame when it has no
    result there: material that left the field of view, that moved into an area that the region
    of interest of a later frame excludes (a crack, a hole), or that fell outside the area
    measured from a reference frame in incremental tracking. With the option on, such points are left
    transparent. With it off, the \ui{Reference frame} view fills them with the value of the
    nearest tracked point inside the region of interest. On the \ui{Deformed frame} view an
    untracked point has no position to be drawn at, so it stays empty either way.
    The choice is also used for exported images and animations; exported data files always
    keep NaN at untracked points.
\end{description}

\subsection{VISUALIZATION}

\begin{description}[leftmargin=1.2cm,style=nextline]
  \item[Colormap]
    The list offers ten colormaps: \file{jet}, \file{viridis}, \file{turbo}, \file{coolwarm},
    \file{plasma}, \file{inferno}, \file{magma}, \file{RdBu\_r}, \file{seismic} and
    \file{black\_rainbow}. The default is
    \file{turbo}, and \file{magma} for the residual. \file{coolwarm}, \file{RdBu\_r} and
    \file{seismic} are diverging maps: with them, the automatic range is symmetric about zero,
    so white means zero. Choose one of them to see the sign of $u$ or $v$ at a glance.
    \file{black\_rainbow} is \file{jet} with its green band replaced by black. The black sits at
    the middle of the colour range, not at zero, and the automatic range is not symmetric with
    this map: give it a \ui{Fixed} range symmetric about zero if black should mean zero.
  \item[Range]
    \ui{Auto} (the default) sets the colour range to the 2nd and 98th percentiles of the values
    drawn on the current frame, after the unit conversion, and computes it again for every frame.
    The extreme 2\,\% at each end take the end colours. \ui{Fixed} keeps \ui{Min} and \ui{Max}
    for every frame. When you switch to \ui{Fixed}, \ui{Min} and \ui{Max} start from the range
    on screen. The boxes hold three decimals, from $-1000$ to $1000$.
  \item[Opacity]
    The opacity of the field over the image, from 0 (transparent) to 100 (opaque); default 70.
\end{description}

The colormap and the range are remembered for each field: switching from \ui{Disp U} to
\ui{Disp V} restores what you last set for \ui{Disp V}. The strain window shares these settings
for its \ui{U}, \ui{V} and \ui{Magnitude} fields (Section~\ref{sec:strain}).

\begin{tipbox}
In \ui{Auto}, each frame has its own range, so one colour means different values on different
frames. Use \ui{Fixed} when you compare frames, play the sequence, or make figures and
animations.
\end{tipbox}

\subsection{PHYSICAL UNITS}
\label{sec:physical-units}

\begin{description}[leftmargin=1.2cm,style=nextline]
  \item[Use physical units]
    Off by default: displacements are in pixels and velocity in pixels per frame. The two
    fields below are enabled only while it is on.
  \item[Pixel size]
    The physical size of one image pixel, with its unit: nm, \textmu m, mm, cm, m or inch,
    followed by \ui{/ px}. Default 1.0\,mm/px. The box holds four decimals: pick the unit in which
    your pixel size has enough digits (5.12\,\textmu m rather than 0.0051\,mm).
  \item[Frame rate]
    The acquisition rate of the images, in frames per second (\ui{fps}); default 1.000\,fps.
\end{description}

A line below the fields shows the resulting units, for example \emph{Disp: mm Velocity: mm/s},
or \emph{Disp: px Velocity: px/fr} while the option is off. With the option on:

\begin{itemize}
  \item \ui{Disp U}, \ui{Disp V} and \ui{Magnitude} are multiplied by the pixel size and shown
        in its unit, on the colorbar, in the cursor read-out and in \ui{Min} and \ui{Max};
  \item the \ui{Velocity} field of the strain window (Section~\ref{sec:strain-fields}) is
        multiplied by the pixel size and the frame rate and shown in the length unit per second;
  \item the \ui{Analysis} tab can plot against time (Section~\ref{sec:analysis});
  \item the residual stays in grey levels, strains stay dimensionless and angles stay in
        degrees.
\end{itemize}

The setting changes what is shown, not the results: data files keep displacements in pixels and
record the pixel size and the frame rate (Section~\ref{sec:export}). The main window and the
strain window share one setting; a change in either shows in both.

\begin{notebox}
A \ui{Fixed} range keeps its numbers when you switch units. After you turn \ui{Use physical
units} on or off, or change the pixel size, set \ui{Min} and \ui{Max} again or go back to
\ui{Auto}.
\end{notebox}

\subsection{The colorbar}

The colorbar sits at the right edge of the canvas. Its title names the field and its unit
(\emph{U~(px)}, \emph{V~(mm)}, \emph{Residual (grey levels)}), and five ticks run from \ui{Min}
at the bottom to \ui{Max} at the top. It follows the range of the frame on show and disappears
when no field is drawn: on a frame that was not computed (after \ui{Cancel}), and while you edit
a region of interest. The cursor read-out and the frame navigator are described in
Section~\ref{sec:canvas}.

\subsection{LOG}

The \ui{LOG} at the bottom of the right sidebar lists the messages of the application, each with
the time (hours, minutes and seconds). It scrolls to the newest line by itself. Errors are red
and completed steps green. % CHECK: warnings sent with level "warning" (engine, Stopped., OOM retry) show in the info colour, only level "warn" is amber; see report.
A run writes its start, the model, the tiling, the frames with a low coverage, the
tracking-quality note and its end there (Section~\ref{sec:running}); an error writes its full
traceback there. \ui{Clear} empties the \ui{LOG}; it changes nothing else. The strain window has
a \ui{LOG} of its own.

\section{Strain}
\label{sec:strain}

Strain is computed from the displacement results in a window of its own, \ui{Strain
Post-Processing}, called the strain window in this guide (Figure~\ref{fig:strain-window}). It
opens by itself when a run finishes; \ui{Open Strain Window} in the right sidebar opens it at
any time once there are displacement results (without them, the \ui{LOG} says \emph{Run DIC
first -{}- no displacement results to post-process.}). The window has two tabs: \ui{Strain
Field}, described here, and \ui{Analysis}, which places probes, gauges and areas on the results
and plots them (Section~\ref{sec:analysis}).

\screenshot{fig-strain-window.png}{The strain window, \ui{Strain Field} tab, showing
\ui{εyy} around the two holes of the plate in tension. Left: the toolbar, the canvas with
its colorbar and the frame navigator. Right: \ui{STRAIN PARAMETERS}, the buttons, \ui{FIELD} and
\ui{VISUALIZATION}.}{fig:strain-window}

\subsection{The strain window}

The left part repeats the canvas of the main window: a toolbar with the fit, \ui{100\%}, zoom-in
and zoom-out buttons, the canvas with its colorbar, and a frame navigator with the same buttons and
playback speeds as the main one (Section~\ref{sec:canvas}). The toolbar also says why nothing is
drawn when that happens (\emph{Click Compute Strain to show this field.}, \emph{Frame $k$ was not
computed.}), and shows the position and value under the pointer.

The right column is a scrollable stack of sections; click a header to fold or unfold it.
\ui{STRAIN PARAMETERS}, \ui{FIELD} and \ui{VISUALIZATION} start unfolded, \ui{ADVANCED},
\ui{PHYSICAL UNITS} and \ui{LOG} folded. The buttons \ui{Compute Strain}, \ui{Cancel} and
\ui{Export Results} sit between the parameters and \ui{FIELD}, outside any section, so they stay
in reach when the sections are folded.

The strain window is independent of the main window. It has its own frame, so moving its slider
does not move the main canvas, and its own \ui{Show on}, background, opacity and edge options.
Two things are shared: \ui{PHYSICAL UNITS} (Section~\ref{sec:physical-units}), and the colormap
and range of \ui{U}, \ui{V} and \ui{Magnitude}, which are those of \ui{Disp U}, \ui{Disp V} and
\ui{Magnitude} in the main window.

\begin{tipbox}
The window opens on frame~1, the reference, where every field is zero. Move its slider to a later
frame to see the deformation.
\end{tipbox}

Closing the window keeps the strain result; \ui{Open Strain Window} brings it back as you left
it. Closing it during a strain computation cancels the computation.

\subsection{How strain is computed}

Strain comes from a virtual strain gauge (VSG). Around each point, RAFTcorr takes a square window
of $L_g \times L_g$ pixels, the gauge, and fits a polynomial by least squares to each displacement
component $u$ and $v$ over the valid points in it. The slopes of the fit at the centre are the
displacement gradients $\partial u/\partial x$, $\partial u/\partial y$, $\partial v/\partial x$
and $\partial v/\partial y$. They give the deformation gradient
\[
  \mathbf{F} = \mathbf{I} + \nabla\mathbf{u}, \qquad
  \mathbf{F} = \begin{pmatrix} 1 + \partial u/\partial x & \partial u/\partial y \\
  \partial v/\partial x & 1 + \partial v/\partial y \end{pmatrix},
\]
from which every strain field is computed. The gradients are taken with respect to the frame-1
positions, so strain is Lagrangian like the displacement. RAFTcorr's displacement is dense (one
value per pixel), and the strain is computed at every pixel of the region of interest unless you
ask for a coarser spacing (\ui{Output spacing}).

A point is left without strain when less than a quarter of its gauge holds valid points
(weighted by the gauge's weighting), when the valid points are arranged so that the fit is
ill-conditioned (a thin strip, for example), and on isolated parts of the region that hold too few
points for a fit. A point near an edge also loses its strain to the trim described below.

\subsection{STRAIN PARAMETERS}

\screenshot[0.5\textwidth]{fig-strain-params.png}{\ui{STRAIN PARAMETERS} with the
\ui{ADVANCED} section unfolded below it.}{fig:strain-params}

The section holds four settings (Figure~\ref{fig:strain-params}); the less used ones are under
\ui{ADVANCED}, below it.

\begin{description}[leftmargin=1.2cm,style=nextline]
  \item[Method]
    The polynomial fitted in each gauge. \ui{Plane fitting} (the default) is the standard
    first-order gauge: a plane, so the strain is the mean gradient over the gauge.
    \ui{Quadratic fitting} adds the second-order terms; it follows a strain gradient inside the
    gauge better, at the cost of more noise. Consider it with a large gauge over a strong
    gradient, such as the edge of a hole.
  \item[VSG size]
    The side $L_g$ of the gauge, from 3 to 401\,px; default 41\,px. It is odd, so that the gauge
    has a centre pixel; an even value moves to the next odd one. A larger gauge averages more
    points: smoother strain, less noise, but a lower spatial resolution, and peaks narrower than
    the gauge are flattened. A smaller gauge resolves sharper features with more noise. The width
    of the trimmed band below grows with it.
  \item[Trim low-confidence edges]
    The coefficient $\alpha$ of the edge trim, from 0.00 to 1.00 in steps of 0.05; default 0.70.
    See below.
  \item[Strain type]
    The strain measure: \ui{Infinitesimal}, \ui{Eulerian} or \ui{Green-Lagrange} (the default).
    See below.
\end{description}

\subsubsection*{Trim low-confidence edges}
\label{sec:edge-trim}

Next to an edge, part of a gauge falls outside the valid region and the fit sees points on one
side only. Such a one-sided fit is biased by the strain gradient, which is often largest at the
edge (a hole, a notch), and it is noisier because fewer points enter it. The trim leaves these
points empty rather than reporting a value that looks as reliable as the rest.

An edge is any border of the valid region of the frame: the boundary of the region of interest,
a hole or a crack, an area that became untracked, and the border of the displacement grid. A point
is trimmed when an edge enters the central square of its gauge whose half-width is
$\alpha\,(L_g - 1)/2$, the coefficient times the half-width of the gauge. With the default gauge
of 41\,px and $\alpha = 0.70$, this half-width is $0.70 \times 20$\,px $= 14$\,px, and the band
along a straight edge is 14\,px wide. Put another way, a point next to a straight edge is kept
when the edge cuts off at most about $(1-\alpha)/2$ of its gauge width, 15\,\% at the default.

\begin{itemize}
  \item $\alpha = 0$: no trim; every point with a fit is reported.
  \item $\alpha = 0.70$: the default.
  \item $\alpha = 1$: only points whose whole gauge lies inside the valid region are reported.
\end{itemize}

The trim is computed for every frame on that frame's valid region, so the band follows a crack
that grows or an area that is lost. Trimmed points are empty (NaN) in every strain field and in
exported data files. Below the box, \emph{Trimmed: $N$ points ($p$\%)} gives how many points of
the frame on show the computed strain lost to the trim. It describes the strain on show, computed
with the value of the last \ui{Compute Strain}, not the value now in the box. The empty band can
be filled in pictures with \ui{Fill trimmed edges (display only)} (Section~\ref{sec:strain-viz}).

\begin{warnbox}
If the band is wider than half the width of the region of interest, every point would be
trimmed, and \ui{Compute Strain} stops with \emph{Edge trimming would remove every strain point
\ldots{} Reduce the VSG size or lower `Trim low-confidence edges'.} Do one of the two, then
compute again.
\end{warnbox}

\subsubsection*{Strain type}

All three measures are computed from $\mathbf{F}$:

\begin{description}[leftmargin=1.2cm,style=nextline]
  \item[Infinitesimal]
    $\boldsymbol{\varepsilon} = \tfrac12\left(\nabla\mathbf{u} + \nabla\mathbf{u}^{T}\right)$.
    Valid for small strains and small rotations only. It is not zero under a rigid rotation: a
    rotation by an angle $\theta$ reads as
    $\varepsilon_{xx} = \varepsilon_{yy} = \cos\theta - 1$.
  \item[Eulerian]
    The Euler-Almansi strain
    $\mathbf{e} = \tfrac12\left(\mathbf{I} - (\mathbf{F}\mathbf{F}^{T})^{-1}\right)$, which
    refers to the deformed configuration. Zero under any rigid rotation.
  \item[Green-Lagrange]
    $\mathbf{E} = \tfrac12\left(\mathbf{F}^{T}\mathbf{F} - \mathbf{I}\right)$, which refers to
    the reference configuration. Zero under any rigid rotation. The default.
\end{description}

For small strains and small rotations the three agree. Keep \ui{Green-Lagrange} when the specimen
rotates or strains are large; choose \ui{Infinitesimal} to compare with a small-strain
calculation, \ui{Eulerian} to compare with a measure in the deformed configuration.

\subsection{ADVANCED}

\begin{description}[leftmargin=1.2cm,style=nextline]
  \item[Weighting]
    How the points of a gauge count in the fit: \ui{Uniform} (the default), all alike, or
    \ui{Gaussian}, decreasing with the distance from the centre as
    $\exp\!\left(-(\Delta x^2 + \Delta y^2)/2\sigma^2\right)$, where $\Delta x$ and $\Delta y$ are
    the offsets from the centre. A Gaussian gauge counts the points near its centre more than
    those near its border.
  \item[Gaussian $\sigma$]
    The $\sigma$ of the Gaussian weight, up to 200\,px; 0 shows \ui{Auto (Lg / 4)}, a quarter
    of the gauge size (10.25\,px for 41\,px). Enabled only with \ui{Gaussian} weighting.
  \item[Output spacing]
    Keep the strain every $n$ pixels in $x$ and $y$, from 1 to 64\,px; default 1\,px (every
    pixel). Each kept point is still fitted over its whole gauge. A larger spacing gives a
    smaller strain result: less memory and smaller exported files for long sequences.
  \item[Pre-fit smoothing]
    Gaussian smoothing of the displacement, with $\sigma$ in pixels up to 50\,px, before the
    strain fit; \ui{Off} (0, the default) differentiates the measured displacement as it is.
    Only the strain uses the smoothed displacement; the displacement results do not change. The
    \ui{Smoothing} parameter of the main window (Section~\ref{sec:parameters}) is different: it
    smooths the displacement results themselves. Smoothing widens the effective gauge.
  \item[Never fit across cracks or holes]
    On by default. Each connected part of the tracked region is fitted on its own, so no gauge
    spans a crack, a hole or the gap between two regions. With it off, a gauge that straddles a
    narrow crack fits one surface over both faces, and the opening of the crack shows as a band
    of high strain as wide as the gauge.
\end{description}

\subsection{Compute Strain}

\ui{Compute Strain} computes the strain of every computed frame (frame~1 included) with the
parameters on show, in the background. A progress bar and a status line appear below the
buttons: \ui{Starting\ldots}, then \ui{Strain frame $i$/$n$}, which counts the frames done, then
\ui{Complete}, which disappears after two seconds. The \ui{LOG} of the window adds
\emph{Strain computation complete.} A project keeps the strain parameters, and the strain itself
when you save the results with it (Section~\ref{sec:projects}).

\ui{Cancel} appears only while the strain is being computed. It stops at the next frame and
keeps the previous strain result (\emph{Strain computation canceled.}). If the computation fails,
a \ui{Strain Computation Failed} dialog gives the reason.

\ui{Export Results} opens the export dialog with the display settings of this window
(Section~\ref{sec:export}). It is available as soon as there are displacement results; the strain
is exported once it is computed.

\subsubsection*{When the strain is out of date}

The strain on show was computed with the parameters of the last \ui{Compute Strain}. As soon as
you change a parameter, a warning line under the buttons reads \ui{Params changed -{}- click
Compute Strain}. Until you click it, the canvas, the \ui{Analysis} tab and exports keep using
the old strain. The line stays until the next \ui{Compute Strain}, even if you set the old values
back.

A new run of the analysis discards the strain, since it belongs to the previous displacements.
The strain buttons then show \emph{Click Compute Strain to show this field.} on the toolbar. A
strain computation still going on when a run finishes is discarded too (\emph{The displacement
results changed during the strain computation; its result was discarded.}).

\subsection{FIELD}
\label{sec:strain-fields}

\ui{FIELD} has two groups of buttons; one button is selected at a time.

Under \ui{DISPLACEMENT}, \ui{U}, \ui{V} and \ui{Magnitude} are the fields of the main window
(Section~\ref{sec:results}). \ui{Velocity} is the speed
$\sqrt{(\partial u/\partial t)^2 + (\partial v/\partial t)^2}$. Its time derivatives are central
differences of the displacement between the neighbouring computed frames, one-sided at the first
and last frame, so it assumes a constant interval between images. It is in pixels per frame, or
in the length unit per second with \ui{Use physical units} on. These four fields need no strain
computation.

Under \ui{STRAIN}, every field comes from the strain tensor of the chosen \ui{Strain type}, or
from $\mathbf{F}$ directly (Table~\ref{tab:strain-fields}). The strains are dimensionless; the two
angles are in degrees. The $y$ axis points down the image, so positive angles turn clockwise on
the screen.

\begin{table}[htbp]
\centering
\small
\caption{The strain fields. $\varepsilon$ stands for the strain tensor of the chosen \ui{Strain
type}.}
\label{tab:strain-fields}
\begin{tabularx}{\textwidth}{@{}l>{\raggedright\arraybackslash}X@{}}
\toprule
Button & Definition \\
\midrule
\ui{εxx}, \ui{εyy} & Normal strains along $x$ (right) and $y$ (down). \\
\ui{εxy} & Shear strain, the tensor component: half the engineering shear,
  $\varepsilon_{xy} = \gamma_{xy}/2$. \\
\ui{ε₁}, \ui{ε₂} & Major and minor principal strains,
  $\tfrac12(\varepsilon_{xx} + \varepsilon_{yy}) \pm
  \sqrt{\tfrac14(\varepsilon_{xx} - \varepsilon_{yy})^2 + \varepsilon_{xy}^2}$,
  with $\varepsilon_1 \geq \varepsilon_2$. \\
\ui{max shear} & Maximum in-plane shear strain, $(\varepsilon_1 - \varepsilon_2)/2$, also a
  tensor value (half the engineering shear). \\
\ui{von Mises} & Equivalent strain
  $\sqrt{\tfrac23(\varepsilon_1^2 + \varepsilon_2^2 + \varepsilon_3^2)}$ with
  $\varepsilon_3 = -(\varepsilon_1 + \varepsilon_2)$, which assumes an incompressible
  material. \\
\ui{ω rot} & Rigid rotation of the material, in degrees, from the polar decomposition of
  $\mathbf{F}$: $\operatorname{atan2}(F_{21} - F_{12},\, F_{11} + F_{22})$. \\
\ui{ε₁ angle} & Direction of $\varepsilon_1$ from $+x$, in degrees between $-90$ and 90:
  $\tfrac12\operatorname{atan2}(2\varepsilon_{xy},\, \varepsilon_{xx} - \varepsilon_{yy})$. \\
\ui{ΔA/A₀} & Relative change of area, $\det\mathbf{F} - 1$. \\
\bottomrule
\end{tabularx}
\end{table}

\subsection{VISUALIZATION}
\label{sec:strain-viz}

The controls work as in the main window (Section~\ref{sec:results}); this subsection lists what
differs.

\begin{description}[leftmargin=1.2cm,style=nextline]
  \item[Show on]
    \ui{Deformed frame} (the default) or \ui{Reference frame}. Strain is computed at the frame-1
    positions; the deformed view draws each value where its point is in the current frame, over
    that frame's image, as the main window does for displacement.
  \item[Background]
    \ui{Show background image}, with \ui{Hidden background} (\ui{White}, \ui{Black},
    \ui{Transparent}) for the hidden case. The window starts from the main window's
    \ui{Hidden background} when it is first opened.
  \item[Colormap, Range, Opacity]
    The same lists and rules as in the main window, remembered for each field. In \ui{Auto} the
    highest 2\,\% of the values, often a peak at the edge of a hole, take the end colour; set a
    \ui{Fixed} range to see the level of the peak. \ui{Min} and \ui{Max} hold six decimals here,
    enough for small strains. \ui{Opacity} is 70 by default.
  \item[Edges]
    \ui{Fill trimmed edges (display only)}, off by default, fills the trimmed band
    (Section~\ref{sec:edge-trim}) with the value of the nearest reliable point of the same
    region. Each region is filled from its own points, so the fill never crosses a crack or a
    hole. It applies to the view and to exported images and animations; exported data files
    always keep the band empty. The filled values are copies, not measurements: leave the option
    off when the edge itself matters.
  \item[Hide untracked regions]
    As in the main window; the window starts from the main window's choice when it is first
    opened, then keeps its own. On the \ui{Reference frame} view, unticking it fills every empty
    point inside the region of interest with the nearest value, the trimmed band included,
    whatever the \ui{Edges} option says. % CHECK: the untracked fill also fills the trimmed band (not region by region); intended? see report.
\end{description}

\ui{PHYSICAL UNITS} is the same group as in the main window and shares its setting. The window's
\ui{LOG} lists the strain computations; failures also open a dialog, since the \ui{LOG} is folded
by default.

\section{Probes and charts: the Analysis tab}
\label{sec:analysis}

The strain window has two tabs. \ui{Strain Field} draws one field over the whole region of
interest (Section~\ref{sec:strain}). \ui{Analysis} shows what happens at one place over the
sequence: the displacement of a point, the mean strain of an area, the reading of a virtual
extensometer or a crack-opening gauge, the field along a line, and, with the record of a testing
machine, a stress--strain curve.

On the left of the tab are the placement tools, the reference image with the probes, and a frame
navigator; on the right, the view tabs, the plot controls, the chart and, below it, the probe
table (Figure~\ref{fig:analysis}). A status line at the bottom repeats the last message, which also
goes to the \ui{LOG}. Probes need a finished run, and strain fields need \ui{Compute Strain} on the
\ui{Strain Field} tab; until then the chart says what is missing.

\screenshot{fig-analysis.png}{The \ui{Analysis} tab with three probes on the reference image: a
virtual extensometer (P1), a point between the holes (P2) and a rectangle (P3). The canvas shows
the plotted field at the current frame. The chart shows each probe's reading in the
\ui{Over time} view, the dashed line marking the current frame; the probe table is below
it.}{fig:analysis}

\subsection{Placing a probe}

The toolbar holds seven placement tools as icons; hover over one to see its name and what it
expects. Click a tool, then click on the image.

\begin{center}
\small
\begin{tabularx}{\linewidth}{@{}l l X@{}}
  \toprule
  Tool & Clicks & Result \\
  \midrule
  \ui{Point}     & one & a point (type \ui{Point} in the table) \\
  \ui{Line}      & two: start and end & a line (type \ui{Line}) \\
  \ui{Rectangle} & two: opposite corners & a region (type \ui{Region}) \\
  \ui{Circle}    & two: centre, then the edge & a region \\
  \ui{Polygon}   & one per vertex, double-click to close & a region of three vertices or more \\
  \ui{Virtual extensometer} & two: the gauge points & a line; the chart plots
    \ui{Extensometer strain} \\
  \ui{Crack gauge} & two: one on each side of the crack & a line; the chart plots
    \ui{Crack opening} \\
  \bottomrule
\end{tabularx}
\end{center}

While a tool is armed, a note on the canvas names it and says what it waits for; \kbd{Esc} cancels
the placement. While you draw a line, its length follows the cursor, in pixels or in the length
unit of the physical units, so you can place an extensometer to a prescribed gauge length $L_0$. A
zero-length line, or a shape without area, is ignored.

The extensometer and the crack gauge are ordinary lines. The two tools only choose what the chart
plots once the line is placed, and only if the chart does not already show a reading of the same
kind. Any line can show any gauge reading.

\begin{notebox}
The canvas always shows the reference image (frame~1), whatever frame the chart is on. Probe
coordinates are reference-image pixels and the fields are stored in reference coordinates
(Section~\ref{sec:conventions}), so a probe follows the material it was placed on through the whole
sequence.
\end{notebox}

\ui{Show field} (on by default) colours the reference image with the plotted field at the current
frame, with the field's own colormap and colour range; for a gauge reading, which has no field, it
shows the field of the \ui{Strain Field} tab. The colorbar uses the unit chosen in
\ui{Strain as}. The buttons after the tools fit the image to the view, show it at \ui{100\%} and
zoom in and out; the mouse wheel zooms as well.

\subsection{Selecting, editing and deleting probes}

With no tool armed:

\begin{itemize}
  \item \textbf{Click} a probe to select it: its row is selected and its curve drawn thicker.
    Click empty space, or press \kbd{Esc}, to clear the selection.
  \item \textbf{Drag} a probe to move it, or drag a square \textbf{handle} of the selected probe to
    reshape it: the ends of a line, the corners of a rectangle, the vertices of a polygon, the
    centre (moves) and the right-hand handle (radius) of a circle. A movement of a few screen
    pixels counts as a click. The chart is recomputed when you release the button.
  \item \textbf{Double-click} a probe, or press \kbd{F2}, to rename it. \kbd{Delete} removes the
    selected probe. There is no undo.
  \item \textbf{Drag empty space}, or drag with the middle mouse button, to pan.
\end{itemize}

Markers and labels keep their size on screen at any zoom; the label of a line also gives its
length. New probes are named \file{P1}, \file{P2}, \ldots\ and take the first free colour of an
eight-colour palette. A deleted probe's number is not reused while other probes remain; once every
probe is deleted, the names start again at \file{P1}.

The probe table has the columns \ui{Show} (untick to hide a probe on the canvas and in the chart
without deleting it), \ui{Name} (double-click to edit; it labels the curve and the CSV columns),
\ui{Type} (\ui{Point}, \ui{Line} or \ui{Region}), \ui{Color} and \ui{Note}, which says why a
probe has gaps or is not plotted. Below it, \ui{Color…} recolours the selected probe, \ui{Delete}
removes it, and \ui{Clear All} removes every probe after a confirmation.

\subsection{Choosing what to plot}

A chart plots one \emph{quantity}, and every visible probe that can produce it is drawn: a point, a
line and a region can share a chart of $\varepsilon_{yy}$, but a strain and a displacement never
share an axis.

\begin{description}[leftmargin=0.6cm,style=nextline]
  \item[Plot]
    A field (\ui{Displacement U}, \ui{Displacement V}, \ui{Displacement magnitude} and the ten
    strain components of Section~\ref{sec:strain}) or, below the separator, a gauge reading of a
    line (Table~\ref{tab:gauge-readings}). Until you choose, the chart follows the field of the
    \ui{Strain Field} tab.
  \item[Statistic]
    How the samples of a line or region become one number per frame: \ui{Mean}, \ui{Median},
    \ui{Maximum}, \ui{Minimum}, \ui{Standard deviation}, or \ui{Valid fraction}, the share of the
    probe's material with a valid value on that frame. A point plots its own value for the first
    four and is left out of the last two. Shown for fields, in the \ui{Over time} and
    \ui{Stress–strain} views.
  \item[Strain as] \ui{ratio}, \ui{\%} or \ui{µε}, for strains (the two angles stay in degrees).
  \item[X axis]
    \ui{Frame} (from 1, as in the application); \ui{Time (s)} when physical units with a frame rate
    are on, frame~1 at 0\,s; \ui{Load (N)} and \ui{Stress (MPa)} with load data (see below).
  \item[Min. valid fraction]
    Default 0.50 (0 to 1). A line or region frame is left blank when fewer than this fraction of its
    points are valid, so that a curve does not stay smooth while its sample shrinks to a few points
    next to a crack. Points are not affected.
\end{description}

A point reads the field by bilinear interpolation. A line is sampled once per grid step along its
length. A region reads every grid point inside it, so its statistics are area-weighted; a region
smaller than one grid cell is read at its centre. Displacements and lengths are in pixels, or in
the length unit when physical units are on; $v$ is positive down, as everywhere in RAFTcorr.

\begin{table}[htbp]
\centering
\small
\caption{Gauge readings of a line. $L_0$ is the drawn length, $L$ the distance between the two
ends on the current frame, $\Delta\mathbf{d}$ the displacement of the second end relative to the
first, $\mathbf{e}$ the unit vector from the first end to the second, and $\mathbf{n}$ that vector
turned to the right on screen.}
\label{tab:gauge-readings}
\begin{tabular}{@{}lll@{}}
  \toprule
  Reading & Value & Meaning \\
  \midrule
  \ui{Extensometer strain}      & $(L - L_0)/L_0$ & engineering strain of the gauge \\
  \ui{Extensometer true strain} & $\ln(L/L_0)$ & logarithmic strain \\
  \ui{Elongation ΔL}            & $L - L_0$ & change of the gauge length \\
  \ui{Crack opening}            & $\Delta\mathbf{d}\cdot\mathbf{e}$ & what a clip gauge reads \\
  \ui{Crack sliding}            & $\Delta\mathbf{d}\cdot\mathbf{n}$ & perpendicular to the gauge,
    along the crack \\
  \ui{Crack opening magnitude}  & $|\Delta\mathbf{d}|$ & the whole separation \\
  \bottomrule
\end{tabular}
\end{table}

A gauge reads the displacement at its two ends only, and reads across a crack by design: an
extensometer that spans a crack includes the opening in its elongation. Place a crack gauge with
its ends on the material on each side of the crack, where they stay tracked in every frame.

\subsection{The four views}

\begin{description}[leftmargin=0.6cm,style=nextline]
  \item[Over time]
    Each probe's reading at every frame; a dashed line marks the current frame. Click the chart to
    go to a frame (with a load or stress axis, the nearest one). The frame is shared with the
    \ui{Strain Field} tab and the navigator under the canvas.
  \item[Along the line]
    The field along the selected line, or the newest visible one, on the current frame, against the
    distance from its first end. With \ui{Other frames} (on by default), up to twelve other frames,
    evenly spaced, are drawn in grey behind it. A gauge reading has no profile: choose a field.
  \item[Kymograph]
    The same line on every frame at once: distance against frame (or time), the value as colour in
    the field's colormap and colour range. Load and stress are not offered as the $x$ axis here,
    because they are not evenly spaced.
  \item[Stress–strain]
    Stress, or load without a cross-section, against the plotted quantity, one curve per probe, the
    current frame ringed. With \ui{Extensometer strain} it is a stress--strain curve; with
    \ui{Crack opening}, a load--opening curve. It needs load data.
\end{description}

The small toolbar above the plot pans, zooms and saves the screen picture; while one of its tools
is active, a click does not change the frame.

\subsection{Gaps and marks}

Every frame of every probe is judged on its own, and a gap in a curve is a gap in the data; it is
never bridged. A frame that has a value but lost part of the probe's material (cut out of that
frame's region of interest, or gone out of the image) gets a hollow marker. In the line views,
lost material has its own grey, because it is neither a zero nor a missing measurement. The
legend and the \ui{Note} column say why a probe has gaps:

\begin{center}
\small
\begin{tabularx}{\linewidth}{@{}l X@{}}
  \toprule
  Note & Meaning \\
  \midrule
  \ui{partly untracked from frame N} & part of the material is lost from frame N on; the value
    comes from what remains \\
  \ui{endpoint lost from frame N} & a gauge end left the tracked material \\
  \ui{gaps: too few valid points} & frames below \ui{Min. valid fraction} \\
  \ui{gaps: no strain value} & the virtual strain gauge found too few points (edge, hole,
    trimmed band): use a smaller gauge or a displacement \\
  \ui{not plotted: off the measured area} & the probe lies outside the region of interest \\
  \ui{not plotted: gauges need a line} & a gauge reading for a point or region \\
  \bottomrule
\end{tabularx}
\end{center}

\subsection{Load data from a testing machine}

\ui{Load data…} imports the machine's record as a CSV file (\file{.csv}, \file{.txt},
\file{.dat}). In the \ui{Load Data} dialog, \ui{Choose file…} reads it and shows its first eight
rows. Comma, semicolon and tab separators, decimal commas, a title line and a units row below the
header (whose units join the column names) are all understood. RAFTcorr guesses the columns from
their names; check them:

\begin{description}[leftmargin=0.6cm,style=nextline]
  \item[Load column] The force, and its unit, \ui{N} or \ui{kN}.
  \item[Match rows to frames: By time]
    A \ui{Time column} and an \ui{Offset}, the machine's time at the reference image (if the camera
    started 2\,s after the machine, enter 2\,s). Frame $n$ is at machine time
    $(n-1)/f + \text{offset}$, with $f$ the frame rate of the physical units; the load is
    interpolated there, and frames outside the record get none. Available only with a frame rate.
  \item[By frame number]
    A \ui{Frame column}, and whether \ui{The first image is numbered} 1 or 0. Rows of the same frame
    are averaged.
  \item[Cross-section A0]
    The initial cross-section in mm$^2$, for the engineering stress $F/A_0$ in MPa; 0 plots the load
    only.
\end{description}

A line at the bottom, \ui{Frames with a load: n of N.}, follows every change, so a wrong offset
shows at once. The record is saved with the project. Loading a new image folder clears it; removing
frames keeps it and warns in the \ui{LOG}. \ui{Remove Load Data}, in the same dialog, clears it by
hand.

\subsection{Exporting probes and charts}

The \ui{Export} menu above the chart is available when the chart shows data:

\begin{description}[leftmargin=0.6cm,style=nextline]
  \item[Probe data (CSV)…]
    In \ui{Over time} and \ui{Stress–strain}. One row per frame: \file{frame} (from 1; frame~1 is the reference
    image), \file{image}, \file{time\_s} with a frame rate, \file{load\_N} and \file{stress\_MPa}
    with load data, then three columns per probe: the value (named after probe, field and
    statistic, as \file{P3\_strain\_eyy\_mean}, or probe and reading, as \file{P1\_strain}), its
    \file{\_valid\_fraction} and a \file{\_flag} (\file{ok}, \file{untracked},
    \file{below\_threshold}, \file{unreliable}, \file{endpoint\_lost}, \file{not\_computed},
    \file{no\_data}). An empty cell means no valid measurement, never 0. Strains are written as
    ratios. Header lines starting with \file{\#} give each probe's geometry in reference pixels, the
    units, the frame numbering, the axis convention, the run and strain settings and the source of
    the load. The file is UTF-8 with a byte-order mark, so spreadsheets read non-Latin names.
  \item[Line data (CSV)…]
    In the line views. One row per sample along the line (distance, \file{x\_px}, \file{y\_px}) and
    one column per frame (\file{frame\_1}, \file{frame\_2}, \ldots).
  \item[Chart image…]
    The chart redrawn for a page: white, 16\,$\times$\,10\,cm, 300\,dpi, PNG, SVG or PDF, text
    editable in SVG and PDF, without the frame cursor and the selection.
  \item[Copy chart, Copy plotted data]
    The same figure, or the numbers as drawn and in the units shown as tab-separated text, to the
    clipboard.
\end{description}

Probes with their names, colours and visibility, and the load record, are saved with the project
(Section~\ref{sec:projects}); the chart settings are not.

\section{Exporting results}
\label{sec:export}

One dialog writes everything RAFTcorr exports: data files, per-frame tables, pictures, animations
and an HTML report. Open it with \ui{Export Results} in the main window, with
\menu{File}{Export Results…} (\kbd{Ctrl+E}), or with \ui{Export Results} in the strain window, once
a run has results. Displacement can be exported before any strain is computed.

The window you open the dialog from sets the starting values of the picture settings:
\ui{Deformed frame} or \ui{Reference frame}, \ui{Show background image}, the hidden-background
fill, \ui{Hide untracked regions}, the opacity and the physical units. From the strain window, the
pictures also follow its \ui{Fill trimmed edges (display only)} (Section~\ref{sec:strain-viz});
from the main window, the trimmed band stays empty. Each field starts with the colormap and colour range
you last gave it in either window. The dialog stays open after an export, so you can write several
formats in turn; \ui{Close} closes it and stops an export still running.

\subsection{Output folder, file names and physical units}

\ui{OUTPUT FOLDER} takes the folder to write to: type it, or pick it with \ui{Browse…}, which also
remembers it for the next time. The export buttons stay disabled until a folder is set; the folder
is created if needed. \ui{Open Folder} opens it after the first successful export.

Every export writes new names, \file{<prefix>\_<kind>\_<timestamp>}: the prefix is the name of the
image folder (spaces and characters Windows does not allow become \file{\_}), the timestamp the
time of the click, \file{YYYYMMDDHHMMSS}. If a file exists already, the dialog lists it and asks
before replacing it.

\begin{center}
\small
\begin{tabularx}{\linewidth}{@{}l X@{}}
  \toprule
  Tab & Written in the output folder \\
  \midrule
  \ui{Data}      & \file{<p>\_parameters\_<t>.json} (always); \file{<p>\_results\_<t>.npz},
                   \file{.mat}, \file{.h5}; \file{<p>\_npz\_<t>/frame\_0001.npz} \ldots;
                   \file{<p>\_csv\_<t>/frame\_0001.csv} \ldots \\
  \ui{Images}    & \file{<p>\_images\_<t>/<field>/frame\_0001.png} \ldots, a folder per field \\
  \ui{Animation} & \file{<p>\_animation\_<t>/<field>.gif} or \file{.mp4}, a file per field \\
  \ui{Report}    & \file{<p>\_report\_<t>.html} \\
  \bottomrule
\end{tabularx}
\end{center}

\file{<field>} is the internal field name, for example \file{disp\_u} or \file{strain\_exx}; hover
over a field in the dialog to see it.

\begin{notebox}
File names, the \ui{FRAME RANGE} boxes and the report count frames from 1, like the frame slider
(\file{frame\_0001} is the reference image); the arrays inside the data files count from 0
(Section~\ref{sec:conventions}). Per-frame files are written for computed frames only, the
reference image included.
\end{notebox}

\ui{PHYSICAL UNITS} holds \ui{Enable physical units}, the \ui{Pixel size} with its unit (\file{nm},
\file{\textmu m}, \file{mm}, \file{cm}, \file{m}, \file{inch} \ui{/ pixel}) and the
\ui{Frame rate} in \ui{fps}, starting from the window's physical units
(Section~\ref{sec:physical-units}).
In pictures, animations and the report, displacements are scaled to the length unit and the
colorbar shows it. Data files stay in pixels: they record the calibration (\file{pixel\_size},
\file{length\_unit}, \file{frame\_rate}), and you multiply by \file{pixel\_size} to convert. Strain
is dimensionless and never scaled. With physical units off, the files record 1, \file{px} and no
frame rate.

\subsection{Data tab}

\screenshot{fig-export-data.png}{The \ui{Data} tab of the export dialog, below the output folder
and the physical units: the formats, the displacement and strain variables, and
\ui{Export Data}.}{fig:export-data}

\ui{FORMAT} chooses the files (Figure~\ref{fig:export-data}). \ui{NumPy Archive (.npz)} and
\ui{MATLAB v7.3 (.mat)} are on by default, \ui{HDF5 (.h5)} and \ui{CSV (per frame)} off. With NPZ
ticked, \ui{NPZ: one file per frame (default: single merged file)} writes one \file{.npz} per
computed frame into a folder of its own. \ui{DISPLACEMENT} and \ui{STRAIN} choose the variables,
with \ui{All} and \ui{None} buttons; \ui{U}, \ui{V}, \ui{Magnitude}, \ui{εxx}, \ui{εyy} and
\ui{εxy} are ticked by default. \ui{Residual} is available when the run recorded one, and the strain
group once you have run \ui{Compute Strain}. NPZ, MAT and HDF5 files always hold the grid
coordinates, the region of interest and the computed-frame flags as well, and the parameters file
is always written.

\subsubsection*{What the data files hold}

NPZ, MAT and HDF5 hold the same variables under the same names (Table~\ref{tab:export-variables}).
$T$ is the number of images, $h \times w$ the bounding box of the region of interest of frame~1
(the grid), and $h_s \times w_s$ the strain grid: the same box, taken every \ui{Output spacing}
pixels of the strain settings (Section~\ref{sec:strain}).

\begin{table}[htbp]
\centering
\small
\caption{Variables of the exported data files. Per-frame arrays are float32, coordinates float64,
flags uint8. The strain variables carry the names of the data files; the strain window shows them
as \ui{εxx}, \ui{εyy}, \ui{εxy}, \ui{ε₁}, \ui{ε₂}, \ui{ε₁ angle}, \ui{max shear}, \ui{von Mises},
\ui{ω rot} and \ui{ΔA/A₀}.}
\label{tab:export-variables}
\begin{tabularx}{\linewidth}{@{}l l l X@{}}
  \toprule
  Name & Shape & Unit & Content \\
  \midrule
  \file{U}, \file{V}        & $(T, h, w)$ & px & displacement, $+u$ right, $+v$ down \\
  \file{Magnitude}          & $(T, h, w)$ & px & $\sqrt{U^2 + V^2}$ \\
  \file{Residual}           & $(T, h, w)$ & grey levels & the residual (when ticked) \\
  \file{X}, \file{Y}        & $(h, w)$ & px & reference-image column and row of each point \\
  \file{ROI}                & $(h, w)$ & & region of interest of frame~1, 1 inside \\
  \file{Done}               & $(T)$ & & 1 where the frame was computed \\
  \file{exx}, \file{eyy}, \file{exy} & $(T, h_s, w_s)$ & & strain components (tensor shear) \\
  \file{e1}, \file{e2}      & $(T, h_s, w_s)$ & & principal strains \\
  \file{principal\_angle}, \file{rotation} & $(T, h_s, w_s)$ & deg & angle of \file{e1} from $+x$;
    rotation \\
  \file{max\_shear}, \file{von\_mises} & $(T, h_s, w_s)$ & & $(\varepsilon_1 - \varepsilon_2)/2$;
    von Mises equivalent strain \\
  \file{area\_change}       & $(T, h_s, w_s)$ & & area change $\Delta A/A_0$ \\
  \file{X\_strain}, \file{Y\_strain} & $(h_s, w_s)$ & px & coordinates of the strain points (with
    any strain variable) \\
  \file{frame\_names}       & $(T)$ & & image file names in frame order \\
  \bottomrule
\end{tabularx}
\end{table}

\begin{itemize}
  \item \textbf{NaN.} Displacement, magnitude and residual are NaN outside the region of interest,
    at points that could not be tracked (cut out by a later frame's region of interest, or gone out
    of the image) and in frames not computed. Strain is NaN outside the region of interest, where
    the virtual strain gauge found too few valid points, in the trimmed band along edges, and in
    frames without strain.
  \item \textbf{Frames and coordinates.} Index 0 is the reference image; its displacement is 0
    inside the region of interest. Every field is in reference (Lagrangian) coordinates. \file{X}
    and \file{Y} are integer pixel positions, the centre of the top-left pixel at $(0, 0)$.
  \item \textbf{Provenance.} Every file carries a JSON document: the RAFTcorr version, creation
    time, frame names, the reference of every frame (\file{ref\_indices}, 0-based), the run
    information (model name and SHA-256 of its file, iterations, correlation mode, device, the run
    settings as run, duration, notes), the grid, the strain settings, the calibration, the
    conventions, and the shape and unit of every variable.
\end{itemize}

In \textbf{HDF5}, each variable is a compressed dataset with a \file{units} attribute; the JSON
document is the root attribute \file{metadata}, beside simple attributes such as \file{n\_frames},
\file{pixel\_size} and \file{frame\_rate} (NaN when unknown). In \textbf{MATLAB}, \file{load} gives
\file{U} as a $T \times h \times w$ array; MATLAB counts from 1, so \file{U(n,:,:)} is frame $n$ of
the slider. The JSON is the text \file{metadata\_json} (read it with \file{jsondecode}), and
\file{frame\_names} has one line per image. In \textbf{NPZ}, the JSON is the string array
\file{metadata\_json}. A \textbf{per-frame NPZ} holds that frame's arrays without the time axis,
the coordinates and the JSON, whose key \file{frame} is the 0-based index.

\begin{lstlisting}[language=Python,showstringspaces=false]
import json, numpy as np
z = np.load("specimen_results_20261119103000.npz")
meta = json.loads(str(z["metadata_json"]))
u = z["U"][4] * meta["pixel_size"]     # frame 5 of the slider, in meta["length_unit"]
\end{lstlisting}

\begin{lstlisting}
S = load('specimen_results_20261119103000.mat');
meta = jsondecode(S.metadata_json);
u5 = squeeze(S.U(5,:,:));               % frame 5 of the slider, in pixels
\end{lstlisting}

The \textbf{parameters file} repeats the provenance as readable JSON. The \textbf{CSV tables} have
one row per point of the region of interest of frame~1: \file{x [px]}, \file{y [px]}, then the
chosen variables with their unit, \file{U [px]}, \ldots, \file{exx [-]}, \ldots,
\file{rotation [deg]}; an empty field is a NaN. With an \ui{Output spacing} above 1\,px the strain
points do not coincide with the displacement rows, so the strain columns go into a second file per
frame, \file{frame\_0001\_strain.csv}.

\begin{warnbox}
A CSV table has one row per pixel of the region of interest: a 1000\,$\times$\,1000\,px region gives
a million rows per frame. Use HDF5, MAT or NPZ for sequences, and CSV for a few frames.
\end{warnbox}

\subsection{Images tab}

\screenshot{fig-export-images.png}{The \ui{Images} tab: one row per field with its colormap, colour
range and opacity, then \ui{IMAGE SETTINGS} and \ui{FRAME RANGE}.}{fig:export-images}

The list at the top has one row per field (Figure~\ref{fig:export-images}): \ui{Export} (ticked by
default for \ui{U}, \ui{V}, \ui{εxx}, \ui{εyy}, \ui{εxy}), the field, its \ui{Colormap}, an
\ui{Auto} or \ui{Fixed} range with \ui{Min} and \ui{Max}, and the \ui{Opacity} ($\alpha$, 0
transparent to 1 opaque). Fields without data are greyed out. \ui{IMAGE SETTINGS}:

\begin{description}[leftmargin=0.6cm,style=nextline]
  \item[Format]
    \ui{PNG} (default), \ui{TIFF}, \ui{JPEG}, \ui{PDF} or \ui{SVG}. PDF and SVG keep the colorbar as
    vectors around an embedded picture of the field.
  \item[Resolution (long edge), JPEG quality, DPI]
    \ui{Full resolution} (default) writes the size set by \ui{DPI} (72 to 600, default 150; the
    field is drawn 7\,inches along its long edge, about 1050\,px at 150\,DPI). A cap of \ui{512 px}
    to \ui{2048 px} limits the long edge of each image and keeps the layout; for PDF and SVG it
    limits the embedded picture. \ui{JPEG quality}: 60 to 100, default 92.
  \item[Auto range]
    For \ui{Auto} rows: \ui{Each frame} (default) takes the range from each frame's own values, as
    the windows do; \ui{Whole sequence} takes one range from all exported frames, so a colour means
    the same value in every image. Use it, or a \ui{Fixed} range, for pictures shown side by side.
  \item[Include colorbar]
    On by default; styled on the \ui{Preview \& Colorbar} tab. It is labelled with the variable name
    of the data files and the unit, for example \file{u (mm)} or \file{exx (-)}.
  \item[Render as]
    \ui{Deformed frame} draws the field where each point is in that frame, over that frame's image;
    \ui{Reference frame} draws it at the positions in frame~1, over frame~1. Untick
    \ui{Show background image} to draw the field alone on the hidden-background fill. \ui{Crop to
    ROI} cuts the picture to the region of interest plus 20 image pixels.
  \item[Hide untracked regions]
    On: points without a result stay transparent. Off: in \ui{Reference frame} pictures they take
    the nearest value inside the region of interest, as on the canvas. Data files keep NaN either
    way.
\end{description}

\ui{FRAME RANGE} exports \ui{All frames}, or \ui{From} \ldots\ \ui{to}, numbered like the slider.
\ui{Export Images} writes one picture per computed frame and ticked field. When the images of the
run cannot be read, the pictures go on white and the \ui{LOG} says so.

\subsection{Animation tab}

The \ui{Animation} tab has its own field list and writes one animation per ticked field.
\ui{Format} is \ui{MP4} (H.264, the default) or \ui{GIF}. MP4 needs the \file{imageio} and
\file{imageio-ffmpeg} packages; without them the MP4 entry is greyed out. A GIF loops and has 255 colours
per frame. \ui{Resolution (long
edge)} defaults to \ui{1024 px}; \ui{Full resolution} draws the images at their own pixel size.
Keep a cap for GIF, whose size grows fast with resolution. \ui{FPS} is 1 to 60 (default 10).
\ui{Frame step} (1 to 20, default 1) keeps every $N$th frame and divides the written frame rate by
$N$, so the animation lasts as long, only choppier. \ui{Include colorbar}, \ui{Render as},
\ui{Crop to ROI} and \ui{Hide untracked regions} work as on the \ui{Images} tab. \ui{Auto} fields
get one colour range for the whole animation. GIF and MP4 have no transparency: a transparent
hidden background is written as white.

\subsection{Report tab}

\ui{Generate Report} writes one HTML file, pictures included, that opens in any browser without a
network connection. Under \ui{CONTENT}, all on by default: \ui{Parameter summary table} (the run and
strain parameters), \ui{Field statistics (min/max/mean/std per frame)} over the valid points of the
region of interest, and \ui{Sample field images} of every $N$th computed frame plus the last
(\ui{Sample every} 5 \ui{frames} by default, 1 to 100). The pictures use each field's row and the
settings of the \ui{Images} tab, with one colour range per field. \ui{FIELDS} chooses the fields;
\ui{U}, \ui{V}, \ui{Magnitude}, \ui{εxx}, \ui{εyy} and \ui{εxy} by default.

\subsection{Preview \& Colorbar tab}

This tab draws one frame with the code of the \ui{Images} export, at a lower resolution: what you
see is what the export writes. Pick a \ui{Field} ticked on the \ui{Images} tab and a \ui{Frame}
(from 1; it starts on frame~2). Transparency shows as a grey checkerboard. \ui{FIELD APPEARANCE}
(\ui{Colormap}, \ui{Range}, \ui{Min}, \ui{Max}, \ui{Opacity}) is the same setting as the field's
row on the \ui{Images} tab; \ui{Apply to all fields} copies the colormap, opacity and range mode,
not \ui{Min} and \ui{Max}, to every ticked field of the \ui{Images} and \ui{Animation} tabs.
\ui{COLORBAR STYLE} applies to every picture, animation and report image:

\begin{center}
\small
\begin{tabularx}{\linewidth}{@{}l l X@{}}
  \toprule
  Setting & Default & Choices \\
  \midrule
  \ui{Position}          & \ui{Right} & \ui{Right}, \ui{Left}, \ui{Top}, \ui{Bottom} \\
  \ui{Font size}, \ui{Font family} & 11\,pt, sans-serif & 6 to 32\,pt; sans-serif, serif,
    monospace \\
  \ui{Bar thickness}     & 0.03 & 0.02 to 0.25 of the picture's width (bar left or right) or
    height (top or bottom) \\
  \ui{Background}        & \ui{White} & \ui{Black} or \ui{White} behind the colorbar and around
    the picture; text takes the other colour \\
  \ui{Margin}, \ui{Margin color} & 0, \ui{White} & 0 to 0.30 of the long edge; \ui{White} or
    \ui{Black} \\
  \ui{Hidden background} & from the window & \ui{White}, \ui{Black} or \ui{Transparent} where the
    image would be; transparency is kept in PNG and TIFF, other formats get white \\
  \bottomrule
\end{tabularx}
\end{center}

\subsection{Running an export}

A running export shows a progress bar and the item being written; \ui{Cancel Export} stops it
after the current frame or file and keeps what was written. Data files are written to a
temporary file and moved into place when complete, so a cancelled or failed export never leaves a
half-written data file and never damages an existing one. Errors appear in red under the button
and in the \ui{LOG}.

\begin{tipbox}
The dialog exports the results it was opened with. After a new run, a new strain computation or
opening a project, it refuses to export and asks you to close it and open it again.
\end{tipbox}

\section{Projects}
\label{sec:projects}

A project file (\file{.raftcorr}) keeps a whole session: which images, the regions of interest, the
settings, the probes and, if you want, the results. Opening it brings the session back without
running the analysis again. The application saves and opens projects; \file{raftcorr run
-{}-project} also writes one, and the Python library reads and writes them
(Section~\ref{sec:cli}).

\subsection{Saving and opening}

\begin{description}[leftmargin=0.6cm,style=nextline]
  \item[\menu{File}{Save Project} (\kbd{Ctrl+S})]
    Saves to the project file of the session. The first time, it asks for a name, starting with
    \file{project.raftcorr} in the image folder; the extension \file{.raftcorr} is added. Images
    have to be loaded first.
  \item[\menu{File}{Save Project As…} (\kbd{Ctrl+Shift+S})]
    Always asks for a name, and the session continues in the new file.
  \item[\menu{File}{Open Project…} (\kbd{Ctrl+O})]
    Opens a project: the images, the regions of interest, the settings, the probes, the load
    record and the results it holds, on the frame that was shown when it was saved.
\end{description}

When the session has results, every save asks \ui{Include the computed results in this project?}
with the approximate uncompressed size of the results. \ui{Yes} (the default) stores the
displacement, the residual and the strain, so the project reopens ready to view, analyse and
export. \ui{No} writes a small file with the settings only, for sharing or as a template; after
opening it you run the analysis again. \ui{Cancel} saves nothing.

Saving and opening run behind a progress dialog. \ui{Cancel} in that dialog stops a save or an
open. A save is written to a temporary file first, so a cancelled or failed save leaves an
existing project as it was. An open reads the stored results frame by frame, the bar follows the
frames read, and a cancelled open leaves the current session as it was. After a save or an open, the title bar shows the name of the project file.

\subsection{What a project holds}

\begin{center}
\small
\begin{tabularx}{\linewidth}{@{}l X@{}}
  \toprule
  Stored & Content \\
  \midrule
  Images & the image folder (its full path and its position relative to the project file) and
    the file names in frame order; not the images themselves \\
  Regions of interest & the region of interest of frame~1 and of every later frame that has one \\
  Settings & tracking mode, model, iterations, smoothing, device, correlation and tiling settings;
    the strain settings; the physical units \\
  View & the displayed field, \ui{Show on}, background, opacity, \ui{Hide untracked regions}, each
    field's colormap and colour range, the current frame, and the last model suggestion \\
  Analysis & the probes with their names, colours and visibility; the load record of a testing
    machine (Section~\ref{sec:analysis}) \\
  Results (optional) & the displacement of every frame, the residual (in half precision), the run
    information, and the strain with its settings and trimmed band \\
  \bottomrule
\end{tabularx}
\end{center}

The images stay where they are. Keep the image folder with the project, or move both together.

\subsection{When the images have moved}

On opening, RAFTcorr looks for the image folder in this order:

\begin{enumerate}
  \item the folder saved in the project;
  \item the folder at the same position relative to the project file, for example after you moved
    or copied the project and the images together, or synced them to another computer. The
    \ui{LOG} says which folder was used;
  \item otherwise it says \ui{The image folder of this project was not found. Choose its new
    location.} and asks for the folder.
\end{enumerate}

The folder must hold every image file of the project under its saved name; if one is missing,
the project does not open and the message names the first missing file. A project saved by a
newer RAFTcorr, with a newer project format, does not open in an older one: the message asks you
to update RAFTcorr.

\subsection{Opening a project by double-click}
\label{sec:projects-doubleclick}

The Windows installer registers the \file{.raftcorr} extension, so a double-click on a project
file starts RAFTcorr and opens it. After a \file{pip} or source installation on Windows, use
\menu{File}{Associate .raftcorr files with RAFTcorr…} once. It registers the extension for your
user account only, needs no administrator rights, and confirms with \emph{``Done.
Double-clicking a .raftcorr file will now open RAFTcorr and restore that project.''}; if the
registry cannot be written, the dialog \ui{File Association Failed} gives the reason. Run it again
if you move or reinstall RAFTcorr, so that the double-click starts the copy you use. The menu
entry exists on Windows only.

On every platform, a project or an image folder can be given on the command line:

\begin{lstlisting}
raftcorr gui path/to/specimen.raftcorr
raftcorr gui path/to/images
\end{lstlisting}

\begin{tipbox}
Save with results before you close RAFTcorr after a long run. A project with results opens
without tracking anything; a settings-only project, or no project, means running the whole
sequence again.
\end{tipbox}

\section{Command line and Python library}
\label{sec:cli}

The application, the \file{raftcorr} command and the Python library share one engine, so the same
images and settings give the same results in all three. The command line suits batches and
computers without a display; the library suits your own scripts. With \file{pip} or a source
installation the command is on the path of that Python environment; the Windows installer puts it
on the path when \ui{Add the ``raftcorr'' command to PATH (new terminals)} is ticked, as it is by
default; on a Mac, use the alias of Section~\ref{sec:installing-macos}.

\subsection{The raftcorr command}

\file{raftcorr} without arguments starts the application. \file{raftcorr -{}-version} prints the
version, \file{raftcorr -h} the commands, and \file{raftcorr <command> -h} the options of one
command.

\begin{center}
\small
\begin{tabularx}{\linewidth}{@{}l X@{}}
  \toprule
  Command & What it does \\
  \midrule
  \file{raftcorr gui [PATH]}      & starts the application, opening an image folder or a
                                    \file{.raftcorr} project if given \\
  \file{raftcorr run FOLDER}      & tracks a sequence and writes a data file and/or a project \\
  \file{raftcorr suggest FOLDER}  & measures the displacement with s8 and suggests a model \\
  \file{raftcorr models}          & lists, downloads or imports the model files \\
  \file{raftcorr selftest}        & runs every installed model on a synthetic image pair \\
  \file{raftcorr info}            & prints the versions of RAFTcorr, Python, PyTorch and the
                                    libraries, and the GPU \\
  \bottomrule
\end{tabularx}
\end{center}

\subsubsection{raftcorr run}

\file{raftcorr run FOLDER [options]} tracks the images of \file{FOLDER} (PNG, JPEG, TIFF, BMP),
sorted in natural order with the first one as the reference image, and writes the results: it is
\ui{Run DIC Analysis} without the window.

\begin{center}
\small
\begin{tabularx}{\linewidth}{@{}l l X@{}}
  \toprule
  Option & Default & Meaning \\
  \midrule
  \file{-{}-roi FILE}    & whole image & region of interest as a mask image, white inside; a mask of
    another size is resized to the images \\
  \file{-{}-model M}     & \file{s4} & \file{s1}, \file{s2}, \file{s4}, \file{s8} or a checkpoint
    path; a missing model file is downloaded \\
  \file{-{}-mode MODE}   & \file{accumulative} & \file{accumulative}, \file{incremental},
    \file{every\_n} or \file{custom} (Section~\ref{sec:workflow}) \\
  \file{-{}-every N}     & 10 & with \file{every\_n}: a new reference every $N$ frames \\
  \file{-{}-keys K,K,…}  & none & with \file{custom}: reference frames as 0-based indices, such as
    \file{0,20,45}; 0 is always one, the last image cannot be one \\
  \file{-{}-iters N}     & 20 & iterations of the model \\
  \file{-{}-device D}    & \file{auto} & \file{auto} (the GPU when CUDA is available), \file{cuda},
    \file{cpu} \\
  \file{-{}-smooth S}    & 0 & Gaussian smoothing of the displacement, $\sigma$ in px; 0 is off \\
  \file{-{}-config FILE} & none & run settings as JSON; replaces \file{-{}-model} to \file{-{}-smooth} \\
  \file{-{}-plain-sort}  & off & sort file names character by character, not naturally \\
  \file{-{}-strain}      & off & also compute the strain \\
  \file{-{}-gauge N}     & 41 & virtual strain gauge size in px, odd \\
  \file{-{}-measure M}   & \file{green\_lagrange} & or \file{infinitesimal}, \file{euler\_almansi} \\
  \file{-{}-edge-trim A} & 0.7 & leave strain empty within $A \times (\text{gauge} - 1)/2$\,px of an
    edge, 0 to 1; 0 keeps every point \\
  \file{-{}-out FILE}    & none & data file: \file{.h5}, \file{.hdf5}, \file{.mat} or \file{.npz} \\
  \file{-{}-project FILE}& none & also save a \file{.raftcorr} project with the results \\
  \file{-{}-overwrite}   & off & replace existing output files \\
  \file{-{}-quiet}       & off & hide the information lines; progress, warnings and errors remain \\
  \bottomrule
\end{tabularx}
\end{center}

The data file holds every variable of Table~\ref{tab:export-variables}, with the residual and, after
\file{-{}-strain}, every strain component. It records no physical units (pixel size 1, unit
\file{px}). The other strain settings keep their defaults: uniform weighting, plane fit, output
spacing 1\,px, no pre-fit smoothing, never across cracks or holes. Regions of interest of later
frames are not available here; use the application or \file{frame\_masks} in Python.

Before tracking, \file{run} checks what would fail at the end: an existing output file without
\file{-{}-overwrite}, a wrong \file{-{}-out} extension, an even gauge, or a trim that would empty the
whole region of interest. Without \file{-{}-out} and \file{-{}-project} it tracks, warns and writes
nothing. Its progress lines, such as \file{[ 25.0\%] Frame 2/5 (reference 1): …}, count frames from
1 like the application.

\begin{lstlisting}
raftcorr run images/ --roi roi.png --model s4 --strain --out results.h5 --project specimen.raftcorr
raftcorr run images/ --model s8 --mode incremental --out results.mat
raftcorr run images/ --mode custom --keys 0,40,80 --out results.npz
raftcorr run images/ --roi roi.png --strain --gauge 31 --edge-trim 0 --out results.h5 --overwrite
raftcorr run images/ --config run.json --out results.h5 --quiet
\end{lstlisting}

A \file{-{}-config} file lists only the keys you set; the others keep their defaults and unknown keys
are ignored. The \file{run\_info.config} entry of an exported parameters file has the same layout,
so you can copy the settings of an earlier run from it.

\begin{lstlisting}
{"model": "s4", "iters": 20, "smoothing_sigma": 0.0,
 "tracking": {"mode": "every_n", "every_n": 20}, "tiling": {"memory_fraction": 0.85}}
\end{lstlisting}

\subsubsection{The other commands}
\label{sec:cli-other}

\begin{description}[leftmargin=0.6cm,style=nextline]
  \item[\file{raftcorr suggest FOLDER [-{}-roi FILE] [-{}-device D]}]
    Measures with s8 the displacement of the region of interest between the reference image and
    up to three later images (a quarter of the way, half way, the last), and prints the finest
    model whose reach is at least 1.3 times the largest displacement (its 99th percentile),
    assuming accumulative tracking. When even s8 falls short, it says to use incremental tracking.
    The frames named in its output are 0-based indices.
  \item[\file{raftcorr models [-{}-download [KEY …]] [-{}-import FILE …]}]
    Lists the four models: name, $s$, reach, the path of the installed file or
    \file{NOT INSTALLED}, a summary, and the folders searched. \file{-{}-download} fetches
    the missing files (all, or the keys given: \file{s1 s2 s4 s8}) into the per-user folder;
    \file{-{}-import} installs \file{.pth} files downloaded by hand from the \file{weights-v1}
    release, recognised by their SHA-256 (Section~\ref{sec:installing-models}).
  \item[\file{raftcorr selftest [-{}-device D] [-{}-gui] [-{}-no-models]}]
    Runs every installed model on a 256\,$\times$\,256\,px synthetic speckle pair with a known
    translation and prints, per model, the measured $u$ and $v$, the time, and \file{ok},
    \file{WRONG}, an error or \file{not installed}. \file{-{}-gui} also opens and closes the main
    window off screen and checks the icons and translations; \file{-{}-no-models} skips the models.
  \item[\file{raftcorr info}]
    Prints the versions of RAFTcorr, Python, the operating system, PyTorch, NumPy, SciPy, PySide6,
    h5py, tifffile, Pillow and Matplotlib, whether CUDA is available, and the GPU with its free
    memory. Add its output to a bug report.
\end{description}

\subsubsection{Exit codes}

\begin{center}
\small
\begin{tabularx}{\linewidth}{@{}c X@{}}
  \toprule
  Code & Meaning \\
  \midrule
  0 & success \\
  1 & \file{selftest}: a model failed, none is installed, or the GUI check failed. Any command: an
      unexpected error, printed as a Python traceback (a folder without images, a reference frame
      out of range, an unknown model key) \\
  2 & a wrong option or value; \file{run}: an output file exists without \file{-{}-overwrite}, a
      wrong \file{-{}-out} extension, or invalid strain settings \\
  3 & \file{models}: a download or an import failed \\
  \bottomrule
\end{tabularx}
\end{center}

\subsection{The Python library}

\begin{lstlisting}[language=Python,showstringspaces=false]
import raftcorr

series = raftcorr.track("images/", roi="roi.png", model="s4")   # DisplacementSeries
strain = raftcorr.compute_strain(series)                        # StrainSeries
u = series.u[4]                        # index 4 = frame 5 on screen; px, +right
exx = strain.components["exx"][4]      # NaN outside the region of interest
\end{lstlisting}

The library follows Section~\ref{sec:conventions}: $+u$ right, $+v$ down, in pixels, in reference
coordinates, and index $k$ is image $k$ counted from 0 (frame $k+1$ in the application).

\begin{lstlisting}[language=Python,showstringspaces=false]
raftcorr.track(images, roi=None, model="s4", config=None, frame_masks=None, callbacks=None)
raftcorr.compute_strain(series, config=None, callbacks=None)
\end{lstlisting}

\begin{center}
\small
\begin{tabularx}{\linewidth}{@{}l X@{}}
  \toprule
  Argument & Meaning \\
  \midrule
  \file{images}      & a folder (natural order), a list of image files, an
    \file{raftcorr.io.ImageSequence}, or a list of 2-D arrays of grey levels 0 to 255 \\
  \file{roi}         & a mask image path (white inside), a boolean array of the image size, or
    \file{None} for the whole image \\
  \file{model}       & \file{s1} to \file{s8} or a checkpoint path, ignored when \file{config} is
    given; a missing model file is downloaded \\
  \file{config}      & \file{track}: a \file{RunConfig}; \file{compute\_strain}: a
    \file{StrainConfig} (default settings of the strain window) \\
  \file{frame\_masks} & \file{\{index: mask\}}, regions of interest of later frames in their own
    coordinates: they exclude what that frame leaves out and are the region of interest when it is
    a reference \\
  \file{callbacks}   & progress, stop and log callbacks (below) \\
  \bottomrule
\end{tabularx}
\end{center}

The settings are frozen dataclasses with the defaults of the application:

\begin{center}
\small
\begin{tabularx}{\linewidth}{@{}l X@{}}
  \toprule
  Class & Fields (default) \\
  \midrule
  \file{RunConfig} & \file{model} (\file{"s4"}), \file{device} (\file{"auto"}), \file{iters} (20),
    \file{corr\_mode} (\file{None} automatic, \file{"tile"} on demand, \file{"volume"} stored),
    \file{smoothing\_sigma} (0.0), \file{tracking}, \file{tiling}; \file{save(path)} and
    \file{load(path)} use the JSON of \file{-{}-config} \\
  \file{TrackingConfig} & \file{mode} (\file{"accumulative"}), \file{every\_n} (10),
    \file{key\_frames} (0-based) \\
  \file{TilingConfig} & \file{margin\_px} (\file{None}: the reach), \file{max\_tile\_pixels}
    (\file{None}: from free GPU memory), \file{memory\_fraction} (0.85), \file{blend\_px} (32) \\
  \file{StrainConfig} & \file{measure} (\file{"green\_lagrange"}), \file{gauge\_px} (41),
    \file{weighting} (\file{"uniform"}, or \file{"gaussian"}), \file{poly\_order} (1, or 2),
    \file{step} (1), \file{gaussian\_sigma} (\file{None}: gauge/4), \file{smoothing\_sigma} (0.0),
    \file{connectivity\_aware} (\file{True}), \file{edge\_trim\_alpha} (0.70); without a control in
    the window: \file{min\_coverage} (0.25), \file{condition\_max} (1000) \\
  \bottomrule
\end{tabularx}
\end{center}

\subsubsection{The result objects}

\begin{center}
\small
\begin{tabularx}{\linewidth}{@{}l X@{}}
  \toprule
  Attribute & Content \\
  \midrule
  \multicolumn{2}{@{}l}{\textbf{DisplacementSeries}} \\
  \file{u}, \file{v}  & $(T, h, w)$ float32, px; NaN outside the region of interest, at untracked
    points and in frames not computed; read-only \\
  \file{mask}, \file{done} & $(h, w)$ region of interest of frame 0; $(T)$ computed frames \\
  \file{grid}         & \file{x0}, \file{y0}, \file{width}, \file{height} and the image size;
    \file{grid.coords()} gives the image coordinates \file{X}, \file{Y} \\
  \file{residual}     & $(T, h, w)$ residual in grey levels, or \file{None} \\
  \file{frame\_names}, \file{n\_frames} & image names; $T$ \\
  \file{completed\_frames} & indices of the computed frames \\
  \file{schedule}     & \file{reference\_of(k)}, the reference of frame $k$ \\
  \file{info}         & \file{model\_name}, \file{model\_sha256}, \file{iters}, \file{device},
    \file{config}, \file{seconds}, \file{notes} \\
  \midrule
  \multicolumn{2}{@{}l}{\textbf{StrainSeries}} \\
  \file{components}   & \file{\{name: (T, h\_s, w\_s)\}} with the names of
    Table~\ref{tab:export-variables}; point $(i, j)$ is at image pixel
    $(x_0 + j\,\text{step},\ y_0 + i\,\text{step})$, with $x_0$, $y_0$ from \file{grid} \\
  \file{step}, \file{done} & output spacing in px; frames with strain \\
  \file{settings}, \file{trimmed} & settings used; edge-trimmed points, or \file{None} \\
  \bottomrule
\end{tabularx}
\end{center}

\subsubsection{Progress and stopping}

A callbacks object has three methods: \file{progress(fraction, message)}, \file{should\_stop()},
polled between frames, and \file{log(message, level)} with the levels \file{debug}, \file{info},
\file{warning} and \file{error}. \file{ConsoleCallbacks} prints like the command line;
\file{FunctionCallbacks} wraps your own functions; the default is silent. When
\file{should\_stop()} returns \file{True}, \file{track} stops after the current frame and returns
what it computed: check \file{series.done}; \file{series.info.notes} records the stop.
\file{compute\_strain} and the exporters raise \file{raftcorr.StopRequested} instead.

\begin{lstlisting}[language=Python,showstringspaces=false]
import raftcorr
from raftcorr import FunctionCallbacks, RunConfig, StrainConfig, TrackingConfig

def show(fraction, message):
    print(f"{100 * fraction:5.1f} %  {message}")

cfg = RunConfig(model="s8", tracking=TrackingConfig(mode="every_n", every_n=20))
series = raftcorr.track("images/", roi="roi.png", config=cfg,
                        callbacks=FunctionCallbacks(progress=show))
strain = raftcorr.compute_strain(series, StrainConfig(gauge_px=31, step=2))
\end{lstlisting}

\subsubsection{Files, probes and projects}

\begin{lstlisting}[language=Python,showstringspaces=false]
from pathlib import Path
from raftcorr.export import PhysicalUnits, export_results, read_results
from raftcorr.analysis import LineGeom, Probe, extract_series
from raftcorr.io.project import Project, save_project

units = PhysicalUnits(scale=0.0125, unit="mm", frame_rate=2.0)   # recorded; data stay in px
export_results("results.h5", series, strain, units=units)        # .h5, .mat or .npz
data = read_results("results.h5")                                # arrays and data["metadata"]

ext = Probe(id=1, kind="line", geometry=LineGeom(100, 200, 400, 200),
            label="E1", color="#ef4444")                          # reference-image pixels
ts = extract_series(series, ext, None, "strain")                  # extensometer strain
ts = extract_series(series, ext, "strain_exx", "mean", strain=strain)

save_project("specimen.raftcorr",
             Project(image_folder=str(Path("images").resolve()),
                     image_files=tuple(series.frame_names), series=series, strain=strain))
\end{lstlisting}

\file{export\_results} writes the data files of the \ui{Data} tab, and \file{read\_results} reads
any of them back with the MATLAB axis order undone. \file{extract\_series} reads a probe as the
\ui{Analysis} tab does and returns \file{frames}, \file{values}, \file{valid\_fraction} and a
\file{status} per frame. It takes the field names of the application (\file{disp\_u},
\file{disp\_v}, \file{disp\_magnitude}, \file{strain\_exx}, \ldots) and, as reduction, a statistic
(\file{value}, \file{mean}, \file{median}, \file{max}, \file{min}, \file{std},
\file{valid\_fraction}) or, for a line, a gauge reading (\file{strain}, \file{true\_strain},
\file{elongation}, \file{cod}, \file{cod\_sliding}, \file{cod\_magnitude}). A project saved this way
opens in the application (Section~\ref{sec:projects}); \file{load\_project} reads one back.

\section{Settings, help and keyboard shortcuts}
\label{sec:settings}

The main window has three menus: \ui{File}, \ui{Settings} and \ui{Help}. This section covers the
\ui{Settings} and \ui{Help} menus, lists the \ui{File} menu with the sections that describe its
items, and gives the keyboard shortcuts. RAFTcorr remembers the theme and the language for the next start, in
the settings of your user account (Section~\ref{sec:overview-files}).

\subsection{The Settings menu}
\label{sec:settings-menu}

\begin{description}[leftmargin=0.6cm,style=nextline]
  \item[\menuthree{Settings}{Theme}{Dark} or \ui{Light}]
    The colour theme of the interface. \ui{Light} is the default. The change is immediate: every
    open window, the strain window included, is drawn again in the new colours, and the log keeps
    its lines.
  \item[\menu{Settings}{Language}]
    The language of the interface: English, Simplified Chinese, Traditional Chinese, Japanese,
    Korean, German, French or Spanish, each listed under its own name. RAFTcorr saves your choice
    and shows \ui{Language changed}, which asks you to restart: the windows already open keep
    their labels until then. The first time RAFTcorr starts, it takes the language of the
    operating system when it has a translation for it (Portuguese falls back to Spanish), and
    English otherwise.
\end{description}

Messages that come from the tracking engine itself, such as the tiling line and the
out-of-memory retries in the \ui{LOG}, stay in English in every language, as does the command
line.

\subsection{The Help menu}
\label{sec:settings-help}

\begin{description}[leftmargin=0.6cm,style=nextline]
  \item[\menu{Help}{Models…}]
    The Models window: which models are installed, downloads, import of model files for an
    offline computer, and the model folder (Section~\ref{sec:installing-models}). It also opens by
    itself when a run needs a model that is missing and cannot be downloaded.
  \item[\menu{Help}{Self-test}]
    Runs every installed model on a synthetic image pair with a known translation and reports
    whether each one finds it (Section~\ref{sec:installing-check}). Use it after an installation,
    an upgrade or a change of PyTorch.
  \item[\menu{Help}{About RAFTcorr}]
    The version number, the four models and the displacement convention: reference to deformed,
    $+u$ right, $+v$ down, in pixels, in the coordinates of frame~1. Give this version number when
    you report a problem (Section~\ref{sec:troubleshooting-report}).
\end{description}

\subsection{The File menu}
\label{sec:settings-file}

\begin{description}[leftmargin=0.6cm,style=nextline]
  \item[\menu{File}{Open Image Folder…}]
    Loads a folder of images, as the drop zone of the \ui{IMAGES} section does
    (Section~\ref{sec:images}).
  \item[\menu{File}{Open Project…}, \ui{Save Project}, \ui{Save Project As…}]
    Projects, \file{.raftcorr} files that hold the images, the regions of interest, the settings
    and optionally the results (Section~\ref{sec:projects}).
  \item[\menu{File}{Export Results…}]
    The export dialog, as the \ui{Export Results} button opens it (Section~\ref{sec:export}). It
    does nothing while there are no results.
  \item[\menu{File}{Associate .raftcorr files with RAFTcorr…}]
    Windows only: makes a double-click on a project file open this copy of RAFTcorr
    (Section~\ref{sec:projects-doubleclick}).
  \item[\menu{File}{Quit}]
    Closes RAFTcorr, as closing the main window does. A run in progress is asked to stop first,
    the strain window closes, and the temporary folder of large results is deleted.
\end{description}

\begin{warnbox}
RAFTcorr does not ask before it closes. Regions of interest, results and probes that are not in a
saved project are lost: save a project first (Section~\ref{sec:projects}).
\end{warnbox}

\subsection{Keyboard shortcuts}
\label{sec:settings-keys}

Table~\ref{tab:shortcuts} lists the shortcuts of the menus and the other keys the application
answers to. On macOS, press \kbd{Cmd} where the table says \kbd{Ctrl}.

\begin{table}[htbp]
\centering
\small
\caption{Keyboard shortcuts.}
\label{tab:shortcuts}
\begin{tabularx}{\linewidth}{@{}>{\raggedright\arraybackslash}p{0.3\linewidth}Xl@{}}
  \toprule
  Keys & Action & Section \\
  \midrule
  \kbd{Ctrl+I} & \menu{File}{Open Image Folder…} & \ref{sec:images} \\
  \kbd{Ctrl+O} & \menu{File}{Open Project…} & \ref{sec:projects} \\
  \kbd{Ctrl+S} & \menu{File}{Save Project} & \ref{sec:projects} \\
  \kbd{Ctrl+Shift+S} & \menu{File}{Save Project As…} & \ref{sec:projects} \\
  \kbd{Ctrl+E} & \menu{File}{Export Results…} & \ref{sec:export} \\
  \kbd{Cmd+Q} & \menu{File}{Quit} on macOS (\kbd{Ctrl+Q} on Linux). Windows has no shortcut
    for it; close the window with \kbd{Alt+F4} & \\
  \midrule
  \kbd{Up}, \kbd{Down}, \kbd{Page Up}, \kbd{Page Down}, \kbd{Home}, \kbd{End}
    & In the image list: move to another frame & \ref{sec:images} \\
  \kbd{Esc} & On the canvas: cancel the shape being drawn & \ref{sec:roi} \\
  \kbd{Esc} & In the \ui{Analysis} tab: cancel the placement of a probe, or deselect & \ref{sec:analysis} \\
  \kbd{Delete} & In the \ui{Analysis} tab: delete the selected probe & \ref{sec:analysis} \\
  \kbd{F2} & In the \ui{Analysis} tab: rename the selected probe & \ref{sec:analysis} \\
  \bottomrule
\end{tabularx}
\end{table}

\section{Troubleshooting}
\label{sec:troubleshooting}

Each entry starts with what you see, then says why it happens and what to do. Look at the
\ui{LOG} panel first: RAFTcorr writes every warning and error there, and the full traceback of a
failed run. Messages quoted in English below appear translated in the other languages, except the
lines that come from the tracking engine itself.

\subsection{Installing}

\begin{description}[leftmargin=0.6cm,style=nextline]
  \item[Windows stops the Setup program: ``Windows protected your PC'']
    The Setup program is not code-signed, and SmartScreen asks about files that few people have
    downloaded. Click \ui{More info}, then \ui{Run anyway} (Section~\ref{sec:installing-windows}).
  \item[macOS refuses to open RAFTcorr]
    The application is not notarised by Apple. Open \ui{System Settings}, \ui{Privacy \&
    Security}, and click \ui{Open Anyway} on the line about RAFTcorr
    (Section~\ref{sec:installing-macos}).
  \item[``RAFTcorr could not install its Python components (code~N)'']
    The second half of the Windows setup failed: PyTorch, RAFTcorr and its dependencies are
    downloaded at that point, so the usual cause is a lost connection, or the console window was
    closed. Code~4 means \emph{``PyTorch could not be installed''}, code~5 \emph{``RAFTcorr could
    not be installed''}. Restore the connection and run \ui{Repair RAFTcorr} from the Start menu.
    The details are in \file{install.log} in the installation folder; search it for
    \file{ERROR}. Run \ui{Repair RAFTcorr} as well when the Start menu entry \ui{RAFTcorr} does
    nothing after an interrupted setup.
  \item[pip stops with ``No such file or directory'' and a hint about long paths]
    Windows limits path lengths, and PySide6 and PyTorch install files deep below
    \file{site-packages}. Create the virtual environment in a short folder such as
    \file{C:\textbackslash venvs\textbackslash raftcorr}, or enable long paths in Windows
    (Section~\ref{sec:installing-pip}).
\end{description}

\subsection{GPU, memory and speed}

\begin{description}[leftmargin=0.6cm,style=nextline]
  \item[RAFTcorr runs on the CPU although the computer has an NVIDIA GPU]
    With \ui{Device} set to \ui{Automatic}, RAFTcorr uses the GPU only when PyTorch sees one, and
    falls back to the CPU without a message. The log line written when the model loads ends with
    the device, for example \emph{``loaded RAFTcorr-s4 (\ldots, cpu)''}, and
    \file{raftcorr info} prints \file{CUDA available: False}. The cause is a CPU build of PyTorch,
    or a driver too old for the CUDA build that is installed. With the Windows Setup program,
    update the NVIDIA driver and run \ui{Repair RAFTcorr}, which replaces a CPU build with the
    CUDA build the driver supports; \file{install.log} says which build it chose and why. With pip or a source installation, run \file{pip uninstall torch} and install
    the CUDA build from the PyTorch index (Section~\ref{sec:installing-pip}). A Mac always runs on
    the CPU.
  \item[``Analysis failed'': ``device 'cuda' was requested but CUDA is not available; use 'cpu'\,'']
    \ui{Device} is set to \ui{GPU (CUDA)}, and PyTorch sees no GPU. Choose \ui{Automatic} or
    \ui{CPU} in \ui{PARAMETERS} (Section~\ref{sec:parameters}), or make the GPU visible as in the
    previous entry.
  \item[``Analysis failed'': ``The GPU ran out of memory. \ldots'']
    The GPU runs out of memory even after RAFTcorr has cut the region into tiles and retried three
    times with smaller windows (Section~\ref{sec:tiling}). Close other GPU programs, a second
    RAFTcorr included; choose a model with a larger $s$, which needs less memory per pixel; or
    lower \ui{GPU Memory} in \ui{ADVANCED} for smaller windows.
  \item[``Analysis failed'': ``There is not enough free GPU memory for the smallest window this
    model accepts. \ldots'']
    The message goes on with \emph{``the memory budget allows windows of about N x N px, below the
    model's M px minimum''}: too little memory is free for even one window, so lowering
    \ui{GPU Memory} cannot help. Free memory (close other GPU programs), raise \ui{GPU Memory} if
    it is set low, or run on the CPU.
  \item[The tiling line ends with ``(margin reduced to fit memory)'']
    The windows that fit the memory are too small to keep the full margin, so a large
    displacement may be missed (Section~\ref{sec:tiling}). Free GPU memory or raise
    \ui{GPU Memory}.
  \item[The analysis is slow on the CPU]
    RAFTcorr runs more slowly on the CPU than on a GPU; the Windows setup says so when it finds
    no NVIDIA driver (\emph{``CPU build (RAFTcorr runs, but slowly)''}). On the CPU, what you can
    change: a smaller region of interest, since the windows cover its bounding box grown by the
    margin; a model with a larger
    $s$, which computes faster and resolves less detail; fewer \ui{Iterations}, which shortens
    every frame but leaves the network fewer refinement steps, so compare with a run at the
    default 20 before you rely on it. \ui{REMAINING} in \ui{PROGRESS} tells you how long the run
    still needs.
\end{description}

\subsection{Models}

\begin{description}[leftmargin=0.6cm,style=nextline]
  \item[The Models window opens when you click Run DIC Analysis]
    The model of the run is not on the computer and could not be downloaded. The card at the top
    of the window gives the reason, for example \emph{``no connection to the download server
    (\ldots); check the internet connection''}, \emph{``cannot write to the model folder
    (\ldots)''} or \emph{``the downloaded file was corrupted (SHA-256 mismatch); try again''}. For
    a run, the log holds the full message with every folder searched. Click \ui{Download} once the connection
    works; on a computer without internet, or behind a firewall that blocks the download, fetch the
    file in a browser and use \ui{Import Model Files…} (Section~\ref{sec:installing-models}).
  \item[A model file is in the folder, but RAFTcorr does not use it]
    Its SHA-256 does not match the release, so RAFTcorr skips it; the log says
    \emph{``\ldots\ exists but its SHA-256 does not match the RAFTcorr-s4 release; skipping it''}.
    The file is damaged or incomplete, or it is not a released model. Download the model again,
    which replaces the file, or import the right file.
  \item[A checkpoint is refused: ``The 1.x models (\ldots) are not supported'']
    \file{-{}-model} on the command line, and the argument \file{model} in Python, also take the
    path of a checkpoint file.
    RAFTcorr~2 loads only checkpoints that store their architecture, such as the four released
    models; the models of RAFTcorr~1.x cannot be used.
\end{description}

\subsection{Tracking}

\begin{description}[leftmargin=0.6cm,style=nextline]
  \item[``Cannot start analysis'']
    Fewer than two images are loaded, or frame~1 has no region of interest, or its region of
    interest holds no pixel (an all-black mask, for example), or the \ui{Reference Frames} field
    holds refused text; the dialog says which (Section~\ref{sec:running}). Load the images or define the region of interest of frame~1
    (Sections~\ref{sec:images} and~\ref{sec:roi}).
  \item[``Tracking Quality Warning'': the motion is beyond the reach of the model]
    The residual jumps from one frame on, typically because the displacement from the reference
    frame exceeds the reach of the model; the warning changes no result
    (Section~\ref{sec:tracking-quality} explains the check). Select \ui{Residual} in \ui{FIELD} to
    see where the images do not match, then:
    \begin{itemize}
      \item click \ui{Suggest a Model}, which measures the motion and picks a model that reaches
        it, or choose a model with a larger $s$ (s8 reaches 256\,px);
      \item in accumulative tracking, switch \ui{Tracking Mode} to \ui{Incremental}, so that only
        each step has to stay within the reach; in incremental tracking, use more reference
        frames (Section~\ref{sec:workflow});
      \item if you turned \ui{Use the Model Reach} off in \ui{ADVANCED}, make the \ui{Margin}
        cover the largest displacement of the region of interest.
    \end{itemize}
  \item[Results look wrong at the edges of the region of interest]
    Look at \ui{Residual} near the edge, then check these causes:
    \begin{itemize}
      \item Points that leave the image cannot be matched, and RAFTcorr drops them. The log warns
        \emph{``Frame N: only P\% of the region of interest could be tracked''} when less than
        90\,\% of the region is left. With \ui{Hide untracked regions} ticked these points are
        transparent; unticked, the \ui{Reference frame} view and pictures fill them from the
        nearest value of the same connected part of the region of interest, while data files
        keep them empty.
      \item A point that moves further than the margin falls outside its window and cannot be
        matched (Section~\ref{sec:tiling}): keep \ui{Use the Model Reach}, or set a \ui{Margin}
        larger than the largest displacement, and watch for \emph{``margin reduced to fit
        memory''} in the tiling line.
      \item A region of interest drawn past the specimen tracks whatever the image shows there,
        the background or a grip. Draw it on the specimen only, cut what lies outside, and give
        later frames a mask of their own when a crack or a hole opens (Section~\ref{sec:roi}).
      \item Strain is biased and partly left empty next to every edge; see \emph{Strain is empty
        next to the edges} below.
    \end{itemize}
  \item[Frames are in the wrong order, or an extra frame appears]
    \ui{Natural Sort (1, 2, …, 10)}, on by default, puts \file{img2} before \file{img10}; untick it
    for names whose plain alphabetical order is the right one. Check the order in the image list
    before you draw the region of interest, since a change of the sort that reorders the frames
    starts a new session, after asking you (Section~\ref{sec:images}). Every image file directly
    inside the folder becomes a frame, a mask image included. \ui{Save} in \ui{REGION OF
    INTEREST} proposes a file next to the image folder, not inside it; keep masks out of the
    image folder, or the next time you load the folder the mask becomes a frame, possibly the
    reference image.
\end{description}

\subsection{Strain}

\begin{description}[leftmargin=0.6cm,style=nextline]
  \item[Strain is empty next to the edges]
    \ui{Trim low-confidence edges} leaves a band along every edge of the region of interest, a
    hole or a crack empty, because a gauge cut by the edge gives a one-sided, biased fit. The band
    is 14\,px wide with the defaults (Section~\ref{sec:edge-trim}). Lower the coefficient to keep
    more points (0 keeps every point, with the bias), or choose a smaller \ui{VSG size}, which
    narrows the band. \ui{Fill trimmed edges (display only)} fills the band in the view and in
    exported pictures; data files always keep it empty.
  \item[``Strain Computation Failed'': ``Edge trimming would remove every strain point'']
    The trimmed band is wider than half the width of the region of interest. Reduce \ui{VSG size}
    or lower \ui{Trim low-confidence edges}.
\end{description}

\subsection{Images}

\begin{description}[leftmargin=0.6cm,style=nextline]
  \item[The image list stays empty]
    RAFTcorr loads the files directly inside the folder with the extensions \file{.png},
    \file{.jpg}, \file{.jpeg}, \file{.tif}, \file{.tiff} and \file{.bmp}; it does not look into
    subfolders. Open the folder that holds the images themselves.
  \item[``Cannot open the images'': ``every frame must have the size of frame 1, \ldots'']
    All images of a sequence must have the size of the first one; the message names the first
    image with its size and lists the files that differ, each with its frame number and size.
    Move those files out of the folder.
  \item[Images of different bit depths, or a dark reference image]
    RAFTcorr converts every frame with one intensity scale, taken from the reference image
    (Section~\ref{sec:images}); the log line after loading names it, for example \emph{``51
    image(s), 2048 x 2048 px, 12-bit (stored as uint16)''}. A later frame stored with another bit
    depth is then scaled wrongly, and later pixels brighter than that range are clipped, without
    a message. Store every frame in the same format and bit depth, and keep the exposure constant
    through the sequence.
\end{description}

\subsection{The application crashes}

\begin{description}[leftmargin=0.6cm,style=nextline]
  \item[``RAFTcorr has hit an error'']
    An unexpected error. The dialog says \emph{``An unexpected error occurred. The application may
    not behave correctly from here on, so saving your project and restarting is recommended.''}
    Save the project, if you can, and restart RAFTcorr. \ui{Show Details} holds the traceback,
    and the log shows it after a line starting with \file{CRASH:}. RAFTcorr also writes it to
    \file{RAFTcorr.log} in its log folder (Section~\ref{sec:overview-files}), and the dialog
    then says \emph{``Details were written to \ldots''} with the path of the file.
  \item[``Analysis failed'' with another message]
    The run stopped on an error. Unlike \ui{Cancel}, an error discards the frames of this run;
    the results of an earlier run stay. The dialog gives the message, and the log holds the full
    traceback.
  \item[RAFTcorr does not start, or closes at once]
    Look at the end of \file{RAFTcorr.log} in the log folder (Section~\ref{sec:overview-files}).
    If it says nothing about the failure, start RAFTcorr from a terminal to see the error:
    \file{raftcorr} (with the Windows Setup program, when the PATH task was ticked), or
    \file{/Applications/RAFTcorr.app/}\allowbreak\file{Contents/MacOS/RAFTcorr} on a Mac. The command
    \file{raftcorr~selftest~-{}-gui} checks that the interface libraries, icons and translations
    load. After an interrupted Windows setup, run \ui{Repair RAFTcorr}.
\end{description}

\subsection{Reporting a problem}
\label{sec:troubleshooting-report}

Report problems at \url{https://github.com/zachtong/RAFTcorr/issues}. Include:

\begin{itemize}
  \item the version (\menu{Help}{About RAFTcorr}, or \file{raftcorr -{}-version}), the operating
    system and how you installed RAFTcorr;
  \item the output of \file{raftcorr info} and, for a model problem, of \menu{Help}{Self-test};
  \item the text of the \ui{LOG} (right-click it, \ui{Select All}, then \ui{Copy}) and the
    \ui{Show Details} text of the error dialog;
  \item \file{install.log} for a problem with the Windows setup, and \file{RAFTcorr.log};
  \item the steps that lead to the problem, and the settings, best as a project saved without
    results.
\end{itemize}

\begin{warnbox}
Issues on GitHub are public. Do not attach images, projects or logs that you may not share; the
log and a project name the folders of your images.
\end{warnbox}


\end{document}
